\documentclass[11pt,a4paper]{article}
\usepackage[T1]{fontenc}
\usepackage[utf8]{inputenc}
\usepackage{lmodern,microtype}
\usepackage[a4paper,margin=25mm]{geometry}
\usepackage{amsmath,amssymb,amsthm,mathtools}
\usepackage{booktabs,array,longtable,enumitem}
\usepackage{xcolor}
\usepackage{hyperref}
\hypersetup{colorlinks=true,linkcolor=blue!45!black,citecolor=blue!45!black,urlcolor=blue!45!black,
 pdftitle={Canonical Digamma Polynomials and the Fifth Harmonic Bridge},
 pdfauthor={Research continuation prepared for the ProveIt project}}
\usepackage{bookmark}
\usepackage{fancyhdr}
\pagestyle{fancy}\fancyhf{}
\fancyhead[L]{\small Canonical digamma polynomials}
\fancyhead[R]{\small ProveIt research continuation}
\fancyfoot[C]{\thepage}
\setlength{\headheight}{14pt}
\setlength{\emergencystretch}{1.5em}
\numberwithin{equation}{section}
\newtheorem{theorem}{Theorem}[section]
\newtheorem{proposition}[theorem]{Proposition}
\newtheorem{lemma}[theorem]{Lemma}
\newtheorem{corollary}[theorem]{Corollary}
\newtheorem{conjecture}[theorem]{Conjecture}
\theoremstyle{definition}
\newtheorem{definition}[theorem]{Definition}
\newtheorem{remark}[theorem]{Remark}
\DeclareMathOperator{\Li}{Li}
\DeclareMathOperator{\Res}{Res}
\DeclareMathOperator{\FP}{FP}
\DeclareMathOperator{\CT}{CT}
\DeclareMathOperator{\Log}{Log}
\newcommand{\C}{\mathbb C}
\newcommand{\Q}{\mathbb Q}
\newcommand{\N}{\mathbb N}
\newcommand{\CM}{\mathcal C}
\newcommand{\re}{\operatorname{Re}}
\title{\textbf{Canonical Digamma Polynomials\protect\\and the Fifth Harmonic Bridge}\\[5pt]
\large Laurent coefficients, polygamma primitives,\\and finite reconstruction of even-tail moments}
\author{Research continuation prepared for the ProveIt project}
\date{11 October 2026\\\small Repository snapshot \texttt{b7667ae40a2957eed9a53d8559d53a3ee20dd1c1}}
\begin{document}
\maketitle
\begin{abstract}
We continue the exact-identity programme in the ProveIt manuscript and
its recent incoming reports. The fifth coordinate finite part of the
digamma function is evaluated in terms of the linear Laurent
coefficients of strict multiple zeta functions of depths two, three,
and four, together with one explicitly specified holomorphic mixed
double-zeta derivative. All finite endpoint constants are proved by
finite harmonic algebra and a normalized Mittag--Leffler expansion.
A Gamma quotient provides a meromorphic realization of the classical
Roman harmonic Bell polynomials. Its coefficients have exactly
prescribed strict-harmonic simple-pole residues and satisfy an ordinary
integral identity isolating one height-one Laurent coefficient at every
depth. We prove uniqueness, up to constants, of the scalar Gamma--Appell
normalization that cancels all nondecaying residue defects, and give
explicit polygamma differential corrections, primitives, and polylogarithmic
integral kernels. Independently, an exponential-polynomial argument proves
that an explicit finite prefix reconstructs every centered moment of
two even harmonic tails; equal orders admit an exact determinant and
the first $q$ moments. A concrete orders-$(2,6)$ consequence is an
absolutely convergent rational polygamma sum. Prior results are attributed,
the remaining stronger conjectures are separated from the proved
statements, and a narrow published Stieltjes sign-description error is
corrected. Exact symbolic replays and independent numerical diagnostics
accompany the proofs.
\end{abstract}
\tableofcontents
\clearpage
\section{Scope, exact targets, and conventions}
\label{sec:scope}

The starting point is the manuscript directory and the incoming archive
directory of the ProveIt repository, inspected at the commit printed on
the title page \cite{Repo}. A repository-relative claim of progress is
different from a claim of priority in the mathematical literature.
This article makes the former precise. Classical Gamma quotients,
harmonic Bell polynomials, height-one coefficient series, and
low-weight multiple-zeta reductions are explicitly attributed. The
proofs developed here organize these tools into a fifth-power bridge
and a finite moment reconstruction theorem. The bounded literature
review is not a proof that no equivalent formulation exists elsewhere.

Two especially relevant incoming reports are \emph{Quartic Tails and
Complex Resolvents} and \emph{Barnes Resolvents and Higher Harmonic Identities}
\cite{Quartic,Barnes}. They already establish the quartic digamma bridge,
the all-depth reciprocal-Gamma singular Laurent coefficients, and finite
formulas for two-even-tail moments. Their stated further questions,
rather than results already proved in their bodies, determine our targets.

\begin{center}
\small
\begin{tabular}{@{}>{\raggedright\arraybackslash}p{.22\textwidth}>{\raggedright\arraybackslash}p{.69\textwidth}@{}}
\toprule
Source target & Result and precise boundary \\
\midrule
Quartic report, Q1: fifth digamma power &
Theorem~\ref{q5:main} gives the complete fifth-power finite part, its
mixed coordinate, and all endpoint constants. \\[5pt]
Quartic report, Q2: canonical polynomial &
The scalar residue normalization is proved for every degree, uniquely
up to an additive constant. An exact differential normalization is also
proved. The stronger claim that scalar polynomials remove every
lower-depth arithmetic coordinate remains open. \\[5pt]
Quartic report, Q8: initial-moment generation &
Theorem~\ref{tr:main} proves an explicit finite-prefix bound and rational
reconstruction. For equal orders the first $q$ moments suffice. The
same first-$q$ assertion for all unequal orders remains open. \\[5pt]
Manuscript $S_6$, revised $S_8$ &
No proof is asserted here. Existing numerical residual certificates
retain their original evidentiary status. \\
\bottomrule
\end{tabular}
\end{center}

\subsection{Harmonic and spectral notation}

Put $H_n^{(r)}=\sum_{k=1}^n k^{-r}$ and $H_n=H_n^{(1)}$, with every
empty sum zero. For strict indices, use
\begin{align}
 H_n(s_1,\ldots,s_d)&=
 \sum_{n\ge k_1>\cdots>k_d\ge1}k_1^{-s_1}\cdots k_d^{-s_d},\\
 Z_d(s_1,\ldots,s_d)&=
 \sum_{k_1>\cdots>k_d\ge1}k_1^{-s_1}\cdots k_d^{-s_d}
 \label{eq:strictzeta}
\end{align}
in their convergence regions, followed by analytic continuation where
specified. The elementary harmonic polynomials are
\begin{equation}\label{eq:edef}
 e_j(n)=H_n(\{1\}^j),\qquad e_0(n)=1,
 \qquad \sum_{j\ge0}e_j(n)t^j=\prod_{k=1}^n\left(1+\frac tk\right).
\end{equation}
They must be distinguished from the complete, weakly ordered harmonic
sums, whose generator has factors $(1-t/k)^{-1}$.

The ordinary Stieltjes convention is
\begin{equation}\label{eq:stieltjes}
 \zeta(1+u)=\frac1u+
 \sum_{j\ge0}\frac{(-1)^j\gamma_j}{j!}u^j,
 \qquad \gamma=\gamma_0.
\end{equation}
At depth $d$, define literal regular Laurent coefficients by
\begin{equation}\label{eq:lambdadef}
 \lambda_{d,j}=[u^j]Z_d(1+u,\{1\}^{d-1}),\qquad
 \lambda_d=\lambda_{d,1}.
\end{equation}
The coefficient brackets impose neither a factorial nor an alternating
sign. In particular
\[
 \lambda_1=-\gamma_1,\qquad
 \lambda_2=\eta,\quad\lambda_3=\theta,\quad\lambda_4=\xi.
\]
The ordered germ has a pole of order $d$, not a simple pole for $d>1$.
It is therefore unsafe to identify it with a Hurwitz-zeta germ merely
by calling both sets of coefficients generalized Stieltjes constants.

We use the reciprocal-Gamma coefficients
\begin{equation}\label{eq:cdef}
 \frac1{\Gamma(1+t)}=\sum_{j\ge0}c_jt^j,
 \quad c_0=1,\quad c_1=\gamma,\quad
 c_2=\frac{\gamma^2-\zeta(2)}2,\quad
 c_3=\frac{\gamma^3-3\gamma\zeta(2)+2\zeta(3)}6.
\end{equation}
They are obtained recursively by differentiating the standard
log-Gamma expansion \cite{DLMF}:
\begin{equation}\label{eq:crecurrence}
 n c_n=\gamma c_{n-1}
       +\sum_{r=2}^n(-1)^{r+1}\zeta(r)c_{n-r}.
\end{equation}

\subsection{A finite part is a specified constant coefficient}

For a function meromorphic at zero and smooth on $(0,1]$, throughout
this article
\begin{equation}\label{eq:fpdef}
 \FP\int_0^1f(x)\,dx
 =\CT_{\varepsilon\downarrow0}\int_\varepsilon^1f(x)\,dx.
\end{equation}
The expansion uses powers of the coordinate $\varepsilon$ itself and
$\log\varepsilon$. There is no silently changed cutoff, rescaled
coordinate, or substituted spectral regularization. If
$f(x)=\sum_{k=-M}^{\infty}a_kx^k$ near zero, this definition can be
computed without divergent numerical cancellation as
\begin{equation}\label{eq:fpsubtraction}
 \int_0^1\left(f(x)-\sum_{k=-M}^{-1}a_kx^k\right)dx
       +\sum_{k=-M}^{-2}\frac{a_k}{k+1}.
\end{equation}
The contribution from $a_{-1}/x$ is zero. For a meromorphic $F$,
\begin{equation}\label{eq:fpderivative}
 \FP\int_0^1F'(x)\,dx=F(1)-[x^0]F(x).
\end{equation}
This rule follows directly by differentiating the Laurent series; it
will account for every endpoint term in the fifth-power calculation.

All ordinary integrals below are identified as such. All matrix ranks
in the moment reconstruction are formal rational ranks. Neither
finite-part calculations nor rational elimination require unproved
arithmetic independence of zeta or Stieltjes constants.

\section{A normalized Mellin formula for the height-one germs}
\label{sec:laurent}

This section gives a self-contained normalization of the coefficient
data used later. The reciprocal-Gamma singular coefficients and
height-one series are established objects in the literature
\cite{Young2023,Young2024} and in the incoming Barnes report
\cite{Barnes}. We record the precise integral form so that the fifth
power and the ordinary Gamma integral share a proved convention.

For $v>0$, put
\begin{equation}\label{eq:Vdef}
 B(v)=-\log(1-e^{-v}),\qquad
 V_d(v)=B(v)^d-\boldsymbol1_{(0,1)}(v)(-\log v)^d.
\end{equation}
Near zero, $B(v)=-\log v+O(v)$, and consequently
\[
 V_d(v)=O\bigl(v(1+|\log v|^{d-1})\bigr).
\]
At infinity it decays exponentially. Thus
$\int_0^\infty v^{u-1}V_d(v)\,dv$ is holomorphic for
$\re u>-1$; each derivative on a compact subset is obtained by
inserting the corresponding power of $\log v$.

\begin{proposition}[Mellin normalization]\label{prop:mellin}
For $d\ge1$, initially for $\re u>0$ and then meromorphically for
$\re u>-1$,
\begin{equation}\label{eq:normalizedmellin}
 \boxed{\displaystyle
 \Gamma(1+u)Z_d(1+u,\{1\}^{d-1})
   =\frac1{u^d}+\frac{u}{d!}
       \int_0^\infty v^{u-1}V_d(v)\,dv.}
\end{equation}
In particular the complete principal part and the constant term are
\begin{equation}\label{eq:principal}
 Z_d(1+u,\{1\}^{d-1})
 =\sum_{j=0}^d c_j u^{j-d}+\lambda_du+O(u^2),
\end{equation}
where
\begin{equation}\label{eq:lambda_integral}
 \lambda_d=c_{d+1}+\frac1{d!}\int_0^\infty\frac{V_d(v)}v\,dv.
\end{equation}
\end{proposition}
\begin{proof}
The generating identity
\[
 \Li_{\{1\}^{d-1}}(y)=\frac{(-\log(1-y))^{d-1}}{(d-1)!}
\]
follows either from strict harmonic generating functions or by
repeated integration of $1/(1-y)$. Termwise Mellin integration gives
\[
 \Gamma(1+u)Z_d(1+u,\{1\}^{d-1})
  =\frac1{(d-1)!}\int_0^\infty
       \frac{v^u B(v)^{d-1}}{e^v-1}\,dv.
\]
For $\re u>0$ the integral is absolutely convergent. Since
$B'(v)=-1/(e^v-1)$, integration by parts, with both boundary terms
zero, turns it into
\[
 \frac{u}{d!}\int_0^\infty v^{u-1}B(v)^d\,dv.
\]
The subtracted comparison integral is exactly
$\int_0^1v^{u-1}(-\log v)^d\,dv=d!/u^{d+1}$.
The estimates preceding the proposition justify continuation and all
coefficient extractions. Multiplication by the Taylor series
\eqref{eq:cdef} proves the remaining formulas.
\end{proof}

\begin{corollary}[Every regular coefficient and spectral derivative]
\label{cor:allcoefficients}
Let
\[
 I_{d,r}=\int_0^\infty \frac{(\log v)^r V_d(v)}v\,dv
 \qquad(d\ge1,\ r\ge0).
\]
For every $j\ge1$,
\begin{equation}\label{eq:alllaurent}
 \lambda_{d,j}=c_{d+j}
  +\frac1{d!}\sum_{r=0}^{j-1}
       c_{j-1-r}\frac{I_{d,r}}{r!}.
\end{equation}
If $\Lambda_d(u)=\Gamma(1+u)Z_d(1+u,\{1\}^{d-1})-u^{-d}$,
then
\begin{equation}\label{eq:normalizedderivatives}
 \Lambda_d^{(j)}(0)=\frac{j}{d!}I_{d,j-1}\qquad(j\ge1).
\end{equation}
These are ordinary convergent logarithmic Mellin integrals. At depth
one, $\lambda_{1,j}=(-1)^j\gamma_j/j!$ recovers the usual Stieltjes
derivative convention.
\end{corollary}
\begin{proof}
Expand the holomorphic integral in \eqref{eq:normalizedmellin} into
its Taylor series and multiply by $\sum c_ju^j$. The estimates are
uniform on every closed disc $|u|\le\rho<1$, so the operations are
analytic coefficient extractions rather than formal regularizations.
\end{proof}

Formula~\eqref{eq:alllaurent} is useful for reproducibility: coefficients
of different depths can be evaluated in one fixed scheme, with no
matching of independently chosen subtraction constants. It is also a
direct route to higher spectral derivatives whenever a new digamma
identity produces them.

\section{The Gamma quotient and its exact harmonic residues}
\label{sec:gamma}

The generating function in this section is classical at integer
arguments: it generates the Roman harmonic numbers, or complete
harmonic star sums, discussed by Sesma \cite[Eq.~(49)]{Sesma}.
We use its meromorphic continuation in the harmonic cutoff to produce
a convenient exact differential normalization of the digamma powers.
The all-depth ordinary integral below is proved directly; its
coefficient-series data are related to Young's established formulas,
and its first case is Coffey's integral for $\gamma_1$
\cite[Eq.~(1.19)]{Coffey}.

\begin{definition}\label{def:Cm}
For $x$ away from the nonpositive integers, define
\begin{equation}\label{eq:GammaGenerator}
 G(x,t)=\frac{\Gamma(1-t)\Gamma(x)}{\Gamma(x-t)}
       =\sum_{m\ge0}\CM_m(x)\frac{t^m}{m!}.
\end{equation}
Coefficient extraction is at $t=0$; the resulting functions of $x$
are extended meromorphically.
\end{definition}

Put
\begin{equation}\label{eq:hr}
 h_1(x)=\psi(x)+\gamma,\qquad
 h_r(x)=\zeta(r)-\zeta(r,x)
       =\zeta(r)+\frac{(-1)^{r+1}}{(r-1)!}\psi^{(r-1)}(x)
       \quad(r\ge2).
\end{equation}
The log-Gamma Taylor expansion yields
\begin{equation}\label{eq:logG}
 \log G(x,t)=\sum_{r\ge1}h_r(x)\frac{t^r}{r}.
\end{equation}
Consequently $\CM_m=Y_m(h_1,1!h_2,\ldots,(m-1)!h_m)$,
where $Y_m$ is the complete exponential Bell polynomial. In particular
\begin{align}
 \CM_0&=1,\qquad \CM_1=h_1,\qquad
 \CM_2=h_1^2+h_2,\\
 \CM_3&=h_1^3+3h_1h_2+2h_3,\\
 \CM_4&=h_1^4+6h_1^2h_2+3h_2^2+8h_1h_3+6h_4,\\
 \CM_5&=h_1^5+10h_1^3h_2+15h_1h_2^2+20h_1^2h_3
                +20h_2h_3+30h_1h_4+24h_5.
 \label{eq:Cfirst}
\end{align}
At every degree a finite recurrence suffices:
\begin{equation}\label{eq:Crecurrence}
 \CM_{m+1}=
 \sum_{r=0}^m\binom mr r!h_{r+1}\CM_{m-r}.
\end{equation}
At a positive integer $n$ the generator is
$\prod_{k=1}^{n-1}(1-t/k)^{-1}$, so $\CM_m(n)/m!$ is the
weakly ordered harmonic sum of $m$ unit indices.

\begin{theorem}[Exact poles and removable values]\label{thm:Cpoles}
For $m\ge2$, $\CM_m$ has only simple poles, with
\begin{equation}\label{eq:Cresidue}
 \Res_{x=-n}\CM_m(x)=-m!e_{m-1}(n),\qquad n\ge0.
\end{equation}
The points $0,-1,\ldots,-(m-2)$ are removable. Its first actual pole
is $-(m-1)$ and has residue $-m$. At the removable points,
\begin{equation}\label{eq:Cnegative}
 \boxed{\displaystyle
 \CM_m(-j)=-m!\sum_{k=0}^j e_k(j)\zeta(m-k),
       \qquad 0\le j\le m-2.}
\end{equation}
In particular $\CM_m(0)=-m!\zeta(m)$, and $\CM_m(1)=0$.
\end{theorem}
\begin{proof}
The reciprocal Gamma factor is entire, so the Gamma quotient has
at most simple poles in $x$. The Gamma recurrence gives exactly
\[
 \Res_{x=-n}G(x,t)
 =\frac{(-1)^n\Gamma(1-t)}{n!\Gamma(-n-t)}
 =-t\prod_{k=1}^n\left(1+\frac tk\right).
\]
Taking coefficients proves the residue assertion. Since $e_{m-1}(n)$
vanishes precisely for $n<m-1$ and is positive thereafter, it also
proves the removability and the first-pole statements.

For the values, set $x=-j+\varepsilon$ and
$E_j(t)=\prod_{k=1}^j(1+t/k)$. Expansion of both Gamma factors gives
the constant coefficient in $\varepsilon$ as
\[
 -tE_j(t)\bigl(H_j-\gamma-\psi(-j-t)\bigr).
\]
By the digamma recurrence this equals
\begin{equation}\label{eq:Cnegativeproof}
 tE_j(t)(\gamma+\psi(1-t))
 -tE_j(t)H_j+E_j(t)+tE_j(t)\sum_{k=1}^j\frac1{k+t}.
\end{equation}
All terms after the first are polynomials of degree at most $j+1$.
For the last term this follows by cancelling the corresponding factor
of $E_j$. Finally
$\gamma+\psi(1-t)=-\sum_{r\ge1}\zeta(r+1)t^r$.
At degree $m\ge j+2$ only this series contributes, proving
\eqref{eq:Cnegative}. Since the pole coefficient at that degree is
zero, its finite Laurent value is the ordinary removable value.
The identity at $1$ follows from $G(1,t)=1$.
\end{proof}

\begin{theorem}[A normally convergent partial-fraction expansion]
\label{thm:CML}
For $m\ge2$, away from its actual poles,
\begin{equation}\label{eq:CML}
 \boxed{\displaystyle
 \CM_m(z)=-m!\zeta(m)-m!\sum_{n\ge1}e_{m-1}(n)
       \left(\frac1{n+z}-\frac1n\right).}
\end{equation}
The series converges normally on every compact subset of that domain.
Terms whose coefficient $e_{m-1}(n)$ is zero are omitted, including
at a removable point where their displayed denominator vanishes.
\end{theorem}
\begin{proof}
Since $e_{m-1}(n)=O((1+\log n)^{m-1})$, the subtracted series
converges normally. To determine the entire remainder, integrate
$z\CM_m(w)/(w(w-z))$ on $|w|=N+1/2$. The right-half-plane
digamma asymptotics and their derivatives, combined with reflection
on the left, give $\CM_m(w)=O((1+\log N)^m)$ uniformly on these
circles. The distance to the nonpositive integral poles stays bounded
below. In the reflection formula the cotangent and its fixed
derivatives are bounded there, while the remaining polygammas have
their usual sectorial bounds \cite{DLMF}. The integral tends to zero
like $O((1+\log N)^m/N)$.

Its residues at $w=z$ and $w=0$ are $\CM_m(z)$ and
$-\CM_m(0)$. Each pole $w=-n$ contributes
$-m!e_{m-1}(n)z/(n(n+z))$. The residue theorem and
Theorem~\ref{thm:Cpoles} give \eqref{eq:CML}. Thus matching the
principal parts alone is not being used to discard an arbitrary
entire function.
\end{proof}

\subsection{An ordinary integral isolates one strict-depth coefficient}

\begin{theorem}[The all-depth ordinary-integral bridge]
\label{thm:Cintegral}
For every integer $m\ge2$,
\begin{equation}\label{eq:Cintegral}
 \boxed{\displaystyle
 \int_0^1\CM_m(x)\,dx=m!\bigl(c_m-\lambda_{m-1}\bigr).}
\end{equation}
The integrand extends continuously to both endpoints. There is no
finite-part operation in this identity.
\end{theorem}
\begin{proof}
Use the opposite-sign parameter
$F(x,t)=G(x,-t)=\Gamma(1+t)\Gamma(x)/\Gamma(x+t)$ and define
\[
 \Phi(t)=\int_0^1\{F(x,t)-1+t(\psi(x)+\gamma)\}\,dx.
\]
Its origin residue is $t-t=0$. More precisely, the subtracted
integrand is jointly holomorphic in $t$ near zero and continuous in
$x\in[0,1]$; this follows by subtracting $t/x$ from the Gamma
quotient. Uniform coefficient extraction on $|t|\le\rho<1$ gives
\begin{equation}\label{eq:PhiTaylor}
 \Phi(t)=\sum_{m\ge2}\frac{(-1)^m t^m}{m!}
                    \int_0^1\CM_m(x)\,dx.
\end{equation}

For $\re t>0$, Euler's beta integral and the standard digamma
integral give, after combining their integrands,
\[
 \Phi(t)=t\int_0^\infty
 \left\{\frac{(1-e^{-v})^t-1}{v}
             +\frac{e^{-v}}{1-e^{-v}}\right\}dv-1.
\]
The combination is integrable at both ends. For example the original
integrand before integration in $x$ is
$e^{-xv}((1-e^{-v})^{t-1}-(1-e^{-v})^{-1})
+e^{-v}/(1-e^{-v})$; its joint integral is absolutely convergent
when $\re t>0$, which justifies Fubini.
Integrate the second term from a lower cutoff $\delta$ and use
$\int_\delta^\infty e^{-v}/(1-e^{-v})\,dv
=-\log(1-e^{-\delta})$. Comparing with
$\int_\delta^1(v^t-1)\,dv/v$ yields
\begin{equation}\label{eq:PhiMellin}
 \Phi(t)=t\int_0^\infty
 \frac{(1-e^{-v})^t-1-\boldsymbol1_{(0,1)}(v)(v^t-1)}v\,dv.
\end{equation}
This subtracted integral is holomorphic for $\re t>-1$ locally
about zero. Near zero, $(1-e^{-v})^t-v^t=O(v^{\re t+1})$
uniformly on compact parameter sets; at infinity the numerator
decays exponentially. Thus the equality continues to $|t|<1$.

At degree $m\ge2$, its right side has coefficient
\[
 \frac{(-1)^{m-1}}{(m-1)!}
             \int_0^\infty\frac{V_{m-1}(v)}v\,dv.
\]
Comparison with \eqref{eq:PhiTaylor} gives
$\int_0^1\CM_m=-m I_{m-1,0}$.
Equation~\eqref{eq:lambda_integral} proves the claimed formula.
\end{proof}

The first two instances are
\begin{align}
 \int_0^1\{(\psi+\gamma)^2+\zeta(2)-\psi'\}\,dx
   &=\gamma^2-\zeta(2)+2\gamma_1,\label{eq:C2integral}\\
 \int_0^1\{(\psi+\gamma)^3
      +3(\psi+\gamma)(\zeta(2)-\psi')
      +2\zeta(3)+\psi''\}\,dx
   &=6(c_3-\eta).
 \label{eq:C3integral}
\end{align}
The apparent singularities in each complete integrand cancel; its
individual displayed terms should not be separately integrated as
ordinary improper integrals. The first formula rearranges exactly to
Coffey's known identity cited above.

\subsection{The known Cauchy-series coordinates recovered from the same quotient}

An additional beta representation is useful both for interpretation
and for stable numerical evaluation:
\begin{equation}\label{eq:Cbeta}
 \CM_m(x)=(x-1)\int_0^1(-\log y)^m(1-y)^{x-2}\,dy,
       \qquad m\ge1,\quad\re x>1-m.
\end{equation}
For $\re x>1$ it follows by differentiating
$G(x,t)=(x-1)B(1-t,x-1)$. The integral on the right then continues
it to the stated half-plane: near $y=1$ the exponent is $x+m-2$.
This also independently displays the location of the first pole.

Define positive rational coefficients
\begin{equation}\label{eq:an}
 a_n=\frac1{(n-1)!}\int_0^1(r)_n\,dr,
 \qquad a_1=\frac12,\quad a_2=\frac56,\quad a_3=\frac98,
\end{equation}
where $(r)_n=r(r+1)\cdots(r+n-1)$ is a rising factorial. If
$L=-\log(1-y)$, their generating kernel is
\begin{equation}\label{eq:Kkernel}
 K(y)=\sum_{n\ge1}a_ny^{n-1}
      =\frac{L-y}{(1-y)^2L^2}.
\end{equation}
Indeed integrate $(1-y)^{-r}=\sum (r)_ny^n/n!$ over $r\in[0,1]$
and differentiate the resulting $y/((1-y)L)$.
Integrating \eqref{eq:Cbeta} in $x$, for $m\ge2$, gives
\[
 \int_0^1\CM_m(x)\,dx
   =-\int_0^1(-\log y)^m K(y)\,dy.
\]
The integrals are absolutely convergent; alternatively Tonelli applies
to their nonnegative negatives. The elementary Mellin integral of
$y^{n-1}$ therefore yields
\begin{equation}\label{eq:knowncauchy}
 \boxed{\displaystyle
 \lambda_d=c_{d+1}+\sum_{n\ge1}\frac{a_n}{n^{d+2}},
       \qquad d\ge1.}
\end{equation}
This coefficient series is an equivalent form of the established
Cauchy/Bernoulli-second-kind series in Young's work
\cite{Young2023,Young2024}; it is included as an independent
normalization check and a connection with the meromorphic differential
realization, not as a new coefficient formula.

\section{A scalar Gamma--Appell normalization of digamma residues}
\label{sec:scalar-gamma-normalizer}

The differential polynomials generated by a normalized Gamma quotient
have exactly the strict harmonic residues required for a single-depth
Laurent coefficient.  There is also a canonical scalar polynomial in
the digamma function that has the same residues up to a decaying
correction.  The distinction is useful: scalar powers are the original
finite-part integrands, whereas differential polynomials remove all
unwanted pole orders at once.

For $n\geq0$, let
\begin{equation}
 e_j(n)=\sum_{1\leq i_1<\cdots<i_j\leq n}
             \frac{1}{i_1\cdots i_j},
 \qquad e_0(n)=1,
 \qquad e_j(n)=0\quad(j>n).
 \label{eq:scalar-elementary-definition}
\end{equation}
Thus $e_j(n)$ is a strict finite harmonic sum, and
\begin{equation}
 E_n(t):=\sum_{j\geq0}e_j(n)t^j
   =\prod_{k=1}^{n}\left(1+\frac{t}{k}\right)
   =\frac{\Gamma(n+1+t)}{\Gamma(n+1)\Gamma(1+t)}.
 \label{eq:scalar-elementary-generator}
\end{equation}
Define the monic polynomials
\begin{equation}
 P_m(Y)=m![t^m]\Gamma(1-t)e^{tY},\qquad m\geq0.
 \label{eq:scalar-appell-definition}
\end{equation}
They satisfy $P_0=1$ and $P_m'(Y)=mP_{m-1}(Y)$, so they form an
Appell sequence.  For example,
\begin{equation}
 \begin{aligned}
 P_1(Y)&=Y+\gamma,\\
 P_2(Y)&=(Y+\gamma)^2+\zeta(2),\\
 P_3(Y)&=(Y+\gamma)^3+3\zeta(2)(Y+\gamma)+2\zeta(3).
 \end{aligned}
 \label{eq:scalar-first-polynomials}
\end{equation}
Every coefficient assertion below is at a fixed finite order; no
uniformity as $m\to\infty$ is required.
The generating function also gives an explicit finite recurrence:
\begin{equation}\label{eq:scalar-recurrence}
 P_{m+1}(Y)=(Y+\gamma)P_m(Y)
  +\sum_{r=1}^m\binom mr r!\zeta(r+1)P_{m-r}(Y).
\end{equation}
It follows by differentiating with respect to $t$ and using
$-\psi(1-t)=\gamma+\sum_{r\ge1}\zeta(r+1)t^r$.

\begin{lemma}[The universal cotangent residue]
\label{lem:scalar-cotangent-residue}
For every $t\in\mathbb C$,
\begin{equation}
 \operatorname*{Res}_{v=0}e^{-\pi t\cot(\pi v)}
       =-\frac{\sin(\pi t)}{\pi}.
 \label{eq:scalar-cotangent-residue}
\end{equation}
Here the residue is the coefficient of $v^{-1}$; the isolated
singularity may be essential.
\end{lemma}

\begin{proof}
Put $w=\pi\cot(\pi v)$.  Near $v=0$, this is an invertible change
between the local coordinate $v$ and the local coordinate $1/w$ at
infinity, since $w=v^{-1}+O(v)$.  Moreover,
\begin{equation}
 \frac{dw}{dv}=-(w^2+\pi^2),\qquad
 e^{-\pi t\cot(\pi v)}\,dv
       =-\frac{e^{-tw}}{w^2+\pi^2}\,dw.
 \label{eq:scalar-residue-change-variable}
\end{equation}
Residues of differentials are invariant under this local coordinate
change.  The two finite residues of the differential on the right are
$-e^{-i\pi t}/(2i\pi)$ at $w=i\pi$ and
$e^{i\pi t}/(2i\pi)$ at $w=-i\pi$.  Their sum is
$\sin(\pi t)/\pi$.  Applying the residue theorem on a circle
containing these two points, or equivalently taking the residue at
infinity, gives its negative, as asserted.  This argument applies to
every complex $t$; a pole rather than an essential singularity at
infinity is not needed for that contour calculation.
\end{proof}

\begin{theorem}[Scalar residue normalization and uniqueness]
\label{thm:scalar-gamma-normalization}
For $m\geq1$ and $n\geq0$, set
\begin{equation}
 r_{m,n}=\operatorname*{Res}_{x=-n}P_m(\psi(x)),
 \qquad
 \delta_{m,n}=r_{m,n}+m!e_{m-1}(n).
 \label{eq:scalar-defect-definition}
\end{equation}
Then $\delta_{1,n}=0$, and for each fixed $m\geq2$,
\begin{equation}
 \delta_{m,n}
   =O_m\!\left(\frac{(1+\log n)^{m-2}}{n}\right)
           \qquad(n\to\infty).
 \label{eq:scalar-defect-decay}
\end{equation}
In particular, the difference in
\eqref{eq:scalar-defect-definition} tends to zero.

Conversely, a monic polynomial $Q\in\mathbb C[Y]$ of degree $m$
satisfies
\begin{equation}
 \operatorname*{Res}_{x=-n}Q(\psi(x))+m!e_{m-1}(n)
           \longrightarrow0
 \label{eq:scalar-normalization-condition}
\end{equation}
if and only if $Q-P_m$ is constant.  Consequently
$P_m(Y)-P_m(0)$ is the unique such monic polynomial with zero
constant term.
\end{theorem}

\begin{proof}
Fix a circle $|v|=\rho$ with $0<\rho<1/2$.  The reflection formula
and the uniform large-argument expansion of the digamma function give
\begin{equation}
 \psi(-n+v)=\psi(n+1-v)-\pi\cot(\pi v)
           =L_n+K(v)+\eta_n(v),
 \label{eq:scalar-local-reflection}
\end{equation}
where $L_n=\log n$, $K(v)=-\pi\cot(\pi v)$, and
$\eta_n$ is holomorphic on $|v|\leq\rho$ with
$\sup|\eta_n|=O(n^{-1})$.  Only $n\geq1$ is needed for this
asymptotic argument.

Consider, coefficientwise near $t=0$,
\begin{equation}
 R_n(t)=\operatorname*{Res}_{v=0}e^{tK(v)+t\eta_n(v)}.
 \label{eq:scalar-residue-germ}
\end{equation}
Integration on the fixed circle makes $R_n(t)$ an entire function
of $t$.  By Lemma~\ref{lem:scalar-cotangent-residue},
\begin{equation}
 R_n(t)=-\frac{\sin(\pi t)}{\pi}+\Delta_n(t),
 \qquad
 \Delta_n(t)=\operatorname*{Res}_{v=0}
       e^{tK(v)}\bigl(e^{t\eta_n(v)}-1\bigr).
 \label{eq:scalar-residue-germ-error}
\end{equation}
Each fixed Taylor coefficient of $\Delta_n$ is $O(n^{-1})$:
this follows directly by bounding its contour integral on
$|v|=\rho$ for $t$ on a fixed small circle.  More precisely,
$\Delta_n(0)=0$ and
$\Delta_n'(0)=\operatorname*{Res}_{v=0}\eta_n(v)=0$.
Thus $\Delta_n$ starts with $t^2$, a fact that gives the exponent
$m-2$ in \eqref{eq:scalar-defect-decay}.

Taking coefficients in
\eqref{eq:scalar-appell-definition} and using
\eqref{eq:scalar-local-reflection}--\eqref{eq:scalar-residue-germ-error}
therefore yields, for $m\geq2$,
\begin{equation}
 \begin{aligned}
 r_{m,n}
 &=-m![t^m]\Gamma(1-t)e^{tL_n}
                         \frac{\sin(\pi t)}{\pi}
       +O_m\!\left(\frac{(1+L_n)^{m-2}}{n}\right)\\
 &=-m![t^{m-1}]\frac{e^{tL_n}}{\Gamma(1+t)}
       +O_m\!\left(\frac{(1+L_n)^{m-2}}{n}\right).
 \end{aligned}
 \label{eq:scalar-nondecaying-polynomial}
\end{equation}
The last equality is Euler's reflection identity
$\Gamma(1-t)\sin(\pi t)/\pi=t/\Gamma(1+t)$.

For comparison, uniformly for $t$ in a sufficiently small fixed disk,
\begin{equation}
 \frac{\Gamma(n+1+t)}{\Gamma(n+1)}
    =e^{tL_n}e^{u_n(t)},\qquad
 u_n(0)=0,\qquad [t^j]u_n(t)=O_j(n^{-1})\quad(j\geq1).
 \label{eq:scalar-gamma-ratio-error}
\end{equation}
For example, this follows by integrating
$\psi(n+1+s)-\log n=O(n^{-1})$ from $s=0$ to $s=t$.
Consequently the correction $e^{u_n(t)}-1$ starts with $t$ and
has each fixed coefficient $O(n^{-1})$.  By
\eqref{eq:scalar-elementary-generator}, for $d\geq1$,
\begin{equation}
 e_d(n)=[t^d]\frac{e^{tL_n}}{\Gamma(1+t)}
       +O_d\!\left(\frac{(1+L_n)^{d-1}}{n}\right).
 \label{eq:scalar-elementary-asymptotics}
\end{equation}
Equations \eqref{eq:scalar-nondecaying-polynomial} and
\eqref{eq:scalar-elementary-asymptotics}, with $d=m-1$, prove
\eqref{eq:scalar-defect-decay}.  For $m=1$ the assertion is
immediate from $P_1(\psi)=\psi+\gamma$, whose residue at every
nonpositive integer is $-1$.

It remains to prove necessity and uniqueness.  Let $Q$ satisfy
\eqref{eq:scalar-normalization-condition}.  If $Q-P_m$ is
nonconstant, write its degree as $j\geq1$ and its leading
coefficient as $a\ne0$.  For $j=1$, its residue is exactly
$-a$.  For $j\geq2$, expand the polynomial at
$L_n+K(v)+\eta_n(v)$ in \eqref{eq:scalar-local-reflection}.
The coefficient of $L_n^j$ is constant in $v$, hence contributes
no residue.  The coefficient of $L_n^{j-1}$ has residue
$ja\operatorname*{Res}_{v=0}(K+\eta_n)=-ja$.  All remaining
terms have degree at most $j-2$ in $L_n$ and bounded residue
coefficients.  Hence
\begin{equation}
 \operatorname*{Res}_{x=-n}(Q-P_m)(\psi(x))
      =-ja(\log n)^{j-1}
          +O\bigl((1+\log n)^{j-2}\bigr).
 \label{eq:scalar-uniqueness-leading-term}
\end{equation}
Neither this expression nor the nonzero constant for $j=1$ can
tend to zero, contradicting the two normalization conditions.
Thus $Q-P_m$ is constant.  Conversely, an additive constant has
zero residue and preserves the condition.
\end{proof}

The theorem establishes a uniqueness of functions, not an assertion
about algebraic independence of Euler's constant or zeta values.
It also explains why an unqualified uniqueness statement would be
false: no residue condition can determine an additive constant.

\begin{proposition}[Uniqueness of the complete Appell family]
\label{prop:scalar-appell-family-uniqueness}
Suppose $A(t)\in\mathbb C[[t]]$, $A(0)=1$, and
$Q_m(Y)=m![t^m]A(t)e^{tY}$.  If
\eqref{eq:scalar-normalization-condition} holds for every $m\geq1$,
then $A(t)=\Gamma(1-t)$ as formal power series.
\end{proposition}

\begin{proof}
By Theorem~\ref{thm:scalar-gamma-normalization},
$Q_m-P_m$ is constant for every $m$.  If
$A(t)-\Gamma(1-t)$ were nonzero, let $a_k t^k$, $a_k\ne0$,
be its first nonzero term; necessarily $k\geq1$.  The coefficient
of $t^{k+1}$ in
$(A(t)-\Gamma(1-t))e^{tY}$ contains the nonconstant term
$a_kY$.  This contradicts the conclusion for $m=k+1$.
\end{proof}

\subsection{An explicit finite formula for the defect}

Set $H_0^{(r)}=0$, write $H_n=H_n^{(1)}$, and define
\begin{equation}
 b_n=H_n-\gamma,\qquad
 a_{r,n}=H_n^{(r+1)}+(-1)^{r+1}\zeta(r+1)\quad(r\geq1).
 \label{eq:scalar-local-coefficients}
\end{equation}
The exact convergent local expansion is
\begin{equation}
 \psi(-n+v)=-\frac1v+b_n+\sum_{r\geq1}a_{r,n}v^r,
                         \qquad0<|v|<1.
 \label{eq:scalar-digamma-laurent}
\end{equation}
It follows by applying the recurrence to $\psi(v)$ and then using
the Taylor expansion of $\psi(1+v)$.

\begin{proposition}[Finite harmonic-polynomial defect formula]
\label{prop:scalar-finite-defect}
For every $m\geq2$ and $n\geq0$,
\begin{equation}
 \begin{aligned}
 r_{m,n}=m!\sum_{h=1}^{m}\frac{(-1)^h}{h!}
 [t^{m-h}v^{h-1}]\,
 \Gamma(1-t)
 \exp\left(tb_n+t\sum_{r=1}^{m-2}a_{r,n}v^r\right),\\
 \delta_{m,n}=r_{m,n}+m!e_{m-1}(n).
 \end{aligned}
 \label{eq:scalar-finite-defect-formula}
\end{equation}
The sum over $r$ is empty when $m=2$.  Formula
\eqref{eq:scalar-finite-defect-formula} requires only finitely
many coefficients of the displayed series; it is an exact
polynomial computation, valid also at $n=0$.
\end{proposition}

\begin{proof}
Substitute \eqref{eq:scalar-digamma-laurent} in
\eqref{eq:scalar-appell-definition} and expand
$e^{-t/v}=\sum_{h\geq0}(-t)^h/(h!v^h)$.  The residue from
the term indexed by $h$ is the coefficient of $v^{h-1}$ in
the holomorphic factor.  Only $1\leq h\leq m$ can contribute
to $[t^m]$.  Moreover, an occurrence of $a_{r,n}v^r$ consumes
at least one additional power of $t$, so whenever such a coefficient
occurs, $h\leq m-1$ and $r\leq h-1\leq m-2$.
All higher $a_{r,n}$ can therefore
be discarded before coefficient extraction.  This proves the
first equation; the second is the definition of the defect.
\end{proof}

For a transparent low-degree form, put
\begin{equation}
 A_j(n)=\zeta(j)-H_n^{(j)},\qquad j\geq2.
 \label{eq:scalar-harmonic-tails}
\end{equation}
The first defects are
\begin{equation}
 \begin{aligned}
 \delta_{1,n}&=\delta_{2,n}=\delta_{3,n}=0,\\
 \delta_{4,n}&=-4A_3(n),\\
 \delta_{5,n}&=-20H_nA_3(n)+5A_2(n)^2
                 -20\zeta(2)A_2(n)+25A_4(n).
 \end{aligned}
 \label{eq:scalar-low-degree-defects}
\end{equation}
These formulas can be checked without any infinite summation.
For example, writing $h=H_n$ and $h_j=H_n^{(j)}$, the finite
formula gives
\begin{equation}
 \begin{aligned}
 r_{4,n}&=-4h^3+12hh_2-4h_3-4\zeta(3),\\
 r_{5,n}&=-5h^4+30h^2h_2-20hh_3-20\zeta(3)h
           -10h_2^2+10\zeta(2)h_2\\
        &\hspace{1.5em}+5\zeta(2)^2+5h_4-25\zeta(4).
 \end{aligned}
 \label{eq:scalar-fourth-fifth-residues}
\end{equation}
Subtract the identities
$-24e_3(n)=-4h^3+12hh_2-8h_3$ and
$-120e_4(n)=-5h^4+30h^2h_2-15h_2^2-40hh_3+30h_4$.
The fifth defect then takes the form in
\eqref{eq:scalar-low-degree-defects} after
$\zeta(2)^2=5\zeta(4)/2$ is used.  In particular, scalar
normalization is already exact at the level of residues through
degree three.  At degrees four and five the remaining terms are
explicit decaying harmonic tails, rather than additional
nondecaying lower-depth data.

\subsection{Comparison with exact differential normalization}

In the notation of the Gamma-polygamma construction, let
\begin{equation}
 \mathcal C_m(x)=m![t^m]\frac{\Gamma(1-t)\Gamma(x)}{\Gamma(x-t)}.
 \label{eq:scalar-differential-comparison-definition}
\end{equation}
Taylor expansion of $\log\Gamma(x-t)$ gives the exact
coefficient identity
\begin{equation}
 \sum_{m\geq0}\frac{\mathcal C_m(x)}{m!}t^m
 =\Gamma(1-t)e^{t\psi(x)}
  \exp\left(\sum_{r\geq2}
       \frac{(-1)^{r+1}\psi^{(r-1)}(x)}{r!}t^r\right).
 \label{eq:scalar-differential-comparison-generator}
\end{equation}
Thus the higher derivatives provide an exact correction to the
scalar Appell family.  The first examples are
\begin{equation}
 \begin{aligned}
 \mathcal C_2&=P_2(\psi)-\psi',\\
 \mathcal C_3&=P_3(\psi)-3P_1(\psi)\psi'+\psi'',\\
 \mathcal C_4&=P_4(\psi)-6P_2(\psi)\psi'+4P_1(\psi)\psi''
                  -\psi'''+3(\psi')^2,\\
 \mathcal C_5&=P_5(\psi)-10P_3(\psi)\psi'
                   +10P_2(\psi)\psi''-5P_1(\psi)\psi'''\\
 &\qquad+15P_1(\psi)(\psi')^2-10\psi'\psi''+\psi''''.
 \end{aligned}
 \label{eq:scalar-differential-comparison-examples}
\end{equation}
Indeed the quotient itself has only simple possible poles in $x$,
and its exact residue is
\begin{equation}
 \operatorname*{Res}_{x=-n}
    \frac{\Gamma(1-t)\Gamma(x)}{\Gamma(x-t)}
       =-t\prod_{j=1}^{n}\left(1+\frac{t}{j}\right).
 \label{eq:scalar-differential-exact-residue}
\end{equation}
Coefficient extraction therefore yields
\begin{equation}
 \operatorname*{Res}_{x=-n} \mathcal C_m(x)=-m!e_{m-1}(n),
 \qquad
 \operatorname*{Res}_{x=-n}\bigl(P_m(\psi(x))-\mathcal C_m(x)\bigr)
          =\delta_{m,n}.
 \label{eq:scalar-exact-defect-comparison}
\end{equation}
The scalar theorem cancels all nondecaying residue data.  The
differential family also cancels the explicit defects and every
higher pole order.  These are precise functional statements;
they do not assert that the associated Laurent constants are
arithmetically independent or form a minimal period basis.

\subsection{A complete finite-part formula for inverse-Gamma derivatives}

There is an all-degree form of the cancellation which explains the
fifth-power identity before its raw-power endpoint corrections are
computed. Define
\begin{equation}\label{eq:Adefinition}
 \mathcal A_m(x)=(-1)^m\Gamma(x)\frac{d^m}{dx^m}\frac1{\Gamma(x)},
 \qquad
 \frac{\Gamma(x)}{\Gamma(x-t)}
       =\sum_{m\ge0}\mathcal A_m(x)\frac{t^m}{m!}.
\end{equation}
Then $\mathcal A_1=\psi$, $\mathcal A_2=\psi^2-\psi'$, and
$\mathcal A_{m+1}=\psi\mathcal A_m-\mathcal A_m'$.

\begin{corollary}[All-degree differential finite parts]
\label{cor:Aintegrals}
For $m\ge1$,
\begin{equation}\label{eq:Aintegrals}
 \boxed{\displaystyle
 \FP\int_0^1\mathcal A_m(x)\,dx
 =m![t^m]\frac{\sin(\pi t)}{\pi t}
   -m!\sum_{d=1}^{m-1}(-1)^{m-d-1}c_{m-d-1}\lambda_d.}
\end{equation}
The sum is empty for $m=1$. For odd $m$ the first term is zero;
for $m=2k$ it is $(-1)^k\pi^{2k}/(2k+1)$.
\end{corollary}
\begin{proof}
Multiplying the series for $\mathcal C_m$ by $1/\Gamma(1-t)$ gives
\[
 \mathcal A_m=\sum_{k=0}^m\frac{m!}{k!}
                         (-1)^{m-k}c_{m-k}\mathcal C_k.
\]
Finite parts are linear. The degree-zero integral is one, and
$\FP\int_0^1\mathcal C_1=\gamma$ because
$\int_\varepsilon^1\psi=-\log\Gamma(\varepsilon)$ has constant
term zero. For $k\ge2$ use Theorem~\ref{thm:Cintegral}.
The terms containing only $c_k$ combine to
\[
 m![t^m]\frac1{\Gamma(1-t)\Gamma(1+t)}
     =m![t^m]\frac{\sin(\pi t)}{\pi t},
\]
and the remaining terms are exactly the displayed finite sum.
\end{proof}

In degree five this reads
\begin{equation}\label{eq:A5height}
 \begin{aligned}
 \FP\int_0^1\mathcal A_5(x)\,dx
   ={}&-120\xi+120\gamma\theta
      -60(\gamma^2-\zeta(2))\eta\\
      &+20(\gamma^3-3\gamma\zeta(2)+2\zeta(3))\lambda_1.
 \end{aligned}
\end{equation}
The next section proves exactly how $\mathcal A_5$ differs, after
integration with the specified finite parts, from the scalar power
$\psi^5$. The resulting additional mixed integral is an analytic
quantity, not an omitted integration constant.

% Independent fifth-power contribution. Standard amsmath/amsthm only.
% Labels have prefix q5:. All finite parts use the coordinate x itself.
\section{The fifth digamma power and its mixed coordinate}
\label{q5:section}

This section resolves the fifth-power calculation posed in the inspected
quartic report.  It fixes all endpoint constants and identifies the
additional mixed spectral combination which occurs beside the
height-one Laurent coefficients.  The result does not assert that this
mixed combination is arithmetically independent of the height-one data.

Write $z_j=\zeta(j)$ and $\gamma=\gamma_0$.  Strict multiple zeta
functions have decreasing indices:
\[
 Z_r(s_1,\ldots,s_r)
   =\sum_{n_1>\cdots>n_r\geq1}
        n_1^{-s_1}\cdots n_r^{-s_r}.
\]
For their height-one germs put
\begin{equation}\label{q5:lambda}
 \lambda_r=[u^1]Z_r(1+u,\underbrace{1,\ldots,1}_{r-1}),
 \qquad
 \lambda_1=-\gamma_1,\quad
 \lambda_2=\eta,\quad\lambda_3=\theta,\quad\lambda_4=\xi.
\end{equation}
The bracket extracts the coefficient itself, with no factorial or sign.
The coordinate finite part of an integral is the constant coefficient
of its cutoff expansion in $\varepsilon$ for the integral from
$\varepsilon$ to the stated upper endpoint.

\subsection{The identities}

Define
\begin{align}
 R(s)&=Z_2(s,3)-z_3\zeta(s),\label{q5:R}\\
 J(s)&=2R(s)+2Z_2(s+2,1)+Z_2(s+1,2)
                 -2\gamma\zeta(s+2)+z_2\zeta(s+1).
 \label{q5:J}
\end{align}
Here $R$ is holomorphic in $\operatorname{Re}s>-1$, and $J$ is
holomorphic in $\operatorname{Re}s>0$.  In particular $J'(1)$ is an
ordinary derivative of a holomorphic function, not a coefficient
requiring a new choice of finite-part normalization.

\begin{theorem}[Fifth-power and mixed-polygamma bridges]
\label{q5:main}
Set
\[
 Q_5=\operatorname{FP}\int_0^1\psi(x)^5\,dx,
 \qquad
 I=\operatorname{FP}\int_0^1\psi(x)\psi'(x)^2\,dx.
\]
Then
\begin{equation}\label{q5:Ibridge}
 \boxed{\displaystyle
 I=-2\gamma z_3+\frac{35}{6}z_4+2J'(1).}
\end{equation}
The fifth-power identity is
\begin{equation}\label{q5:bridge}
 \boxed{\begin{aligned}
 Q_5={}&30\gamma^2z_2-60\gamma z_3+\frac{35}{3}z_4
       -120\xi+120\gamma\theta\\
 &-60(\gamma^2-z_2)\eta
       +20(\gamma^3-3\gamma z_2+2z_3)\lambda_1
       +10J'(1).
 \end{aligned}}
\end{equation}
Equivalently, the mixed integral itself isolates the extra coordinate:
\begin{equation}\label{q5:QminusI}
 \boxed{\begin{aligned}
 Q_5-5I={}&30\gamma^2z_2-50\gamma z_3-\frac{35}{2}z_4
       -120\xi+120\gamma\theta\\
 &-60(\gamma^2-z_2)\eta
       +20(\gamma^3-3\gamma z_2+2z_3)\lambda_1.
 \end{aligned}}
\end{equation}
All the ordinary zeta derivatives in the compact notation are explicit:
\begin{align}\label{q5:Jexpanded}
 J'(1)={}&2R'(1)
       +2\left.\frac{\partial}{\partial s}Z_2(s,1)\right|_{s=3}
       +\left.\frac{\partial}{\partial s}Z_2(s,2)\right|_{s=2}
       -2\gamma\zeta'(3)+z_2\zeta'(2).
\end{align}
If $\phi_{13}=[u^1]Z_2(1+u,3)$, then
$R'(1)=\phi_{13}+z_3\gamma_1=\phi_{13}-z_3\lambda_1$.
\end{theorem}

The last conversion matters.  Substituting $\phi_{13}$ for $R'(1)$
without its $z_3\gamma_1$ correction would give an incorrect fifth-power
formula.  The mixed term is one specified linear combination.  No claim
that its separate components form a minimal arithmetic basis is needed.

\subsection{A convergent arithmetic form and the spectral germs}

For $n\geq1$ use
\begin{equation}\label{q5:integer-bcd}
 b=\psi(n),\qquad c=2z_2-\psi'(n),\qquad
 d=\frac12\psi''(n),\qquad L_n=\log n,
\end{equation}
where $b,c,d$ in a summand always refer to that summand's $n$.
Then
\begin{align}
 J(s)&=\sum_{n\geq1}\frac1{n^s}
      \left(2d+\frac{c}{n}+\frac{2b}{n^2}\right),
      &&\operatorname{Re}s>0,\label{q5:Jseries}\\
 I&=-2\gamma z_3+\frac{35}{6}z_4
       -\sum_{n\geq1}
        \left(\frac{4d}{n}+\frac{2c}{n^2}
                         +\frac{4b}{n^3}\right)L_n.
 \label{q5:Iarithmetic}
\end{align}
Another expression for $Q_5$ is
\begin{align}\label{q5:Qarithmetic}
 Q_5={}&30\gamma^2z_2-60\gamma z_3+\frac{35}{3}z_4
           +20\gamma_4-120z_2\gamma_2\notag\\
 &+20\sum_{n\geq1}
     \frac{(b^3-3bc+d-L_n^3+6z_2L_n)L_n}{n}\notag\\
 &-20\sum_{n\geq1}\frac{bL_n}{n^3}
       -10\sum_{n\geq1}\frac{cL_n}{n^2}.
\end{align}
Every series in these displays converges absolutely.  Indeed
$\psi(n)=\log n+O(n^{-1})$,
$\psi'(n)=O(n^{-1})$, and $\psi''(n)=O(n^{-2})$.
The summand in the first sum of \eqref{q5:Qarithmetic} is
$O((1+\log^3n)/n^2)$.  These estimates also prove normal convergence
of \eqref{q5:Jseries} and of each fixed spectral derivative on compact
subsets of its half-plane.  The series for $R$ is
\[
 R(s)=\sum_{n\geq1}
       \frac{H_{n-1}^{(3)}-z_3}{n^s};
\]
its summands are $O(n^{-\operatorname{Re}s-2})$.

The height-one Laurent germs used here have a direct existence proof.
For $r\geq1$ and $\operatorname{Re}s>1$,
\begin{equation}\label{q5:mellin}
 Z_r(s,1^{r-1})=
 \frac{(-1)^{r-1}}{(r-1)!\,\Gamma(s)}
 \int_0^\infty
    \frac{t^{s-1}\log^{r-1}(1-e^{-t})}{e^t-1}\,dt.
\end{equation}
This follows from the elementary-symmetric generating series and
termwise Mellin integration.  Near zero the kernel is its displayed
leading term
$(-1)^{r-1}\log^{r-1}t/((r-1)!t)$ plus
$O(1+|\log t|^{r-1})$.  Subtraction of that term gives
\[
 \Gamma(1+u)Z_r(1+u,1^{r-1})=u^{-r}+V_r(u),
 \qquad V_r(0)=0,
\]
with $V_r$ holomorphic near zero.  The assertion $V_r(0)=0$ follows
by integrating the derivative of $\log^r(1-e^{-t})$ and comparing its
lower boundary with $\log^r t$.  Thus the principal and constant
coefficients are fixed by $1/\Gamma(1+u)$, and in particular $\xi$
in \eqref{q5:lambda} is well defined.

\subsection{Global partial fractions with a fixed entire remainder}

For $n\geq0$, set
\begin{equation}\label{q5:localcoeffs}
 \begin{aligned}
 b_n&=H_n-\gamma,&
 c_n&=z_2+H_n^{(2)},\\
 d_n&=H_n^{(3)}-z_3,&
 e_n&=z_4+H_n^{(4)}.
 \end{aligned}
\end{equation}
The local digamma expansion is
\[
 \psi(-n+\varepsilon)
 =-\varepsilon^{-1}+b_n+c_n\varepsilon
       +d_n\varepsilon^2+e_n\varepsilon^3+O(\varepsilon^4).
\]
For $F_5(z)=\psi(z)^5$ and $F_I(z)=\psi(z)\psi'(z)^2$, let
\[
 P_{X,n}(z)=\sum_{j=1}^5\frac{p_{X,j,n}}{(z+n)^j},
 \qquad X\in\{5,I\}.
\]
The principal coefficients, with subscripts $n$ suppressed in the table,
are
\begin{center}
\begin{tabular}{c|c|c}
$j$&$p_{5,j}$&$p_{I,j}$\\\hline
5&$-1$&$-1$\\
4&$5b$&$b$\\
3&$-10b^2+5c$&$-c$\\
2&$10b^3-20bc+5d$&$2bc-3d$\\
1&$-5b^4+30b^2c-20bd-10c^2+5e$&$4bd+c^2-5e$
\end{tabular}
\end{center}
Write $r_{X,n}=p_{X,1,n}$.  The constant coefficients at zero are
\begin{align}
 \kappa_5={}&-\gamma^5+20\gamma^3z_2-30\gamma^2z_3
         -30\gamma z_2^2+20\gamma z_4+20z_2z_3-5z_5,
 \label{q5:kappa5}\\
 \kappa_I={}&-\gamma z_2^2-6\gamma z_4-2z_2z_3+7z_5.
 \label{q5:kappaI}
\end{align}

\begin{lemma}[Two normalized Mittag--Leffler expansions]
\label{q5:ML}
For $X=5,I$ and $z\notin\{0,-1,-2,\ldots\}$,
\begin{equation}\label{q5:MLformula}
 F_X(z)=P_{X,0}(z)+\kappa_X+
       \sum_{n\geq1}\bigl(P_{X,n}(z)-P_{X,n}(0)\bigr).
\end{equation}
The series converges normally on compact sets avoiding the poles.
\end{lemma}
\begin{proof}
The coefficient table follows by finite multiplication of the local
digamma series and its derivative.  For the entire remainder, put
$f_X=F_X-P_{X,0}$ and integrate
$z f_X(w)/(w(w-z))$ on $|w|=N+1/2$.
Digamma asymptotics in the right half-plane and the reflection formula
in the left half-plane give $f_X(w)=O(1+\log^5N)$ on these circles.
The cotangent and its fixed derivatives are uniformly bounded there:
away from the real axis use their exponential expressions, and on the
remaining arcs the real coordinate stays a fixed distance from the
integers.  The contour integral is therefore
$O((1+\log^5N)/N)$ and tends to zero.
Its residues are $f_X(z)$, $-f_X(0)$, and
$-(P_{X,n}(z)-P_{X,n}(0))$.  This proves the identity.
Finally $p_{5,j,n}=O(1+\log^{5-j}n)$ and the corresponding mixed
coefficients have no greater growth.  The two-point subtractions give
an additional inverse power of $n$, proving normal convergence.
\end{proof}

\subsection{The finite summation ledger}

Integrating \eqref{q5:MLformula} and telescoping each pole of order
at least two gives
\begin{equation}\label{q5:integrated}
 \operatorname{FP}\int_0^1F_X(x)\,dx
   =B_X+\sum_{n\geq1}r_{X,n}
        \left(\log\left(1+\frac1n\right)-\frac1n\right),
\end{equation}
where
\begin{equation}\label{q5:Bmechanism}
 B_X=\kappa_X+
  \sum_{j=2}^5\sum_{n\geq1}
    \left(\frac{p_{X,j,n}-p_{X,j,n-1}}{(j-1)n^{j-1}}
                       -\frac{p_{X,j,n}}{n^j}\right).
\end{equation}
This formula includes the finite parts at the origin.  For example,
the lower endpoint contribution of $p_{X,j,0}x^{-j}$ is
$-p_{X,j,0}/(j-1)$; it cancels the initial term of that pole's finite
telescoping sum.  The simple-pole finite part is zero.

Here is a shorter directly reproducible form of all the nonlogarithmic
terms.  In the sums in this display only, let
$b=b_{n-1}$, $c=c_{n-1}$, $d=d_{n-1}$.  Then
\begin{align}
 B_5={}&\kappa_5+\sum_{n\geq1}\left\{
 \frac{-10b^3+20bc-5d+30b^2-20c}{n^2}
 +\frac{-20b^2+15c}{n^3}
 +\frac{5b-35/6}{n^4}+\frac6{n^5}\right\},
 \label{q5:B5sum}\\
 B_I={}&\kappa_I+\sum_{n\geq1}\left\{
 \frac{-2bc+3d+2c}{n^2}+\frac{2b-c}{n^3}
 +\frac{-3b-7/6}{n^4}+\frac2{n^5}\right\}.
 \label{q5:BIsum}
\end{align}
All these sums converge absolutely.

To evaluate them and the finite residue sums, write
\[
 H_N(a_1,\ldots,a_d)=
    \sum_{N\geq n_1>\cdots>n_d\geq1}
                 n_1^{-a_1}\cdots n_d^{-a_d},
 \qquad e_j(N)=H_N(1^j).
\]
The finite product rule is the quasi-shuffle rule with a positive merged
index.  The summation rule is equally elementary:
\begin{equation}\label{q5:finite-prefix}
 \sum_{n=1}^N\frac{H_n(a_1,\ldots,a_d)}{n^p}
   =H_N(p,a_1,\ldots,a_d)
      +H_N(p+a_1,a_2,\ldots,a_d),
\end{equation}
with the second term absent for an empty word.  Both statements follow
by separating strict inequalities from equalities of the outer indices.
There is no infinite rearrangement in either identity.

For completeness, the convergent values required in this section are
\begin{equation}\label{q5:MZVtable}
\begin{array}{c|c}
\text{value}&\text{reduction}\\\hline
Z_2(2,1)&z_3\\
Z_2(3,1),\ Z_2(2,2),\ Z_3(2,1,1)&z_4/4,\ 3z_4/4,\ z_4\\
Z_2(4,1),\ Z_3(3,1,1)&2z_5-z_2z_3\\
Z_2(3,2),\ Z_3(2,2,1)&3z_2z_3-11z_5/2\\
Z_2(2,3),\ Z_3(2,1,2)&9z_5/2-2z_2z_3\\
Z_4(2,1,1,1)&z_5
\end{array}
\end{equation}
These are the classical low-weight reductions.  Their use here can be
checked from the finite product rule, the depth-two sum theorem and
iterated-integral duality.  At weight five, for example, the stuffle
and shuffle products are respectively
\[
 z_2z_3=Z_2(2,3)+Z_2(3,2)+z_5,
 \qquad
 z_2z_3=Z_2(2,3)+3Z_2(3,2)+6Z_2(4,1).
\]
Together with
$Z_2(4,1)+Z_2(3,2)+Z_2(2,3)=z_5$, these give the three
depth-two weight-five entries.  Substitution $t\mapsto1-t$ in the
iterated integrals gives the paired dual values in the table.
The same arguments at weights three and four give the first two rows;
$z_2^2=5z_4/2$ is used throughout.

Here is an explicit finite identity which records the full fifth-power
constant.  In this display set $h_j=H_N^{(j)}$, $h=h_1$,
$E_j=e_j(N)$, and abbreviate $H_N(a_1,\ldots,a_d)$ by
$H(a_1,\ldots,a_d)$.  Put
\[
 A=\gamma^3-3\gamma z_2+z_3,\qquad
 r_0=-5\gamma^4+30\gamma^2z_2-20\gamma z_3-10z_2^2+5z_4.
\]
Then the exact equality for every $N\geq1$ is
\begin{align}\label{q5:S5finite}
 \sum_{n=1}^N\frac{r_{5,n}}n={}&
  -120E_5+120\gamma E_4+60(z_2-\gamma^2)E_3
       +20AE_2+r_0h+20Ah_2\notag\\
 &-120H(2,1,1,1)+120\gamma H(2,1,1)
       +(50z_2-60\gamma^2)H(2,1)\notag\\
 &+20E_2h_3-20H(3,1,1)-20hh_4
       +20H(4,1)+20h_5\notag\\
 &+10hH(2,2)-10H(2,1,2)-10H(2,2,1)+10H(2,3)\notag\\
 &-20\gamma hh_3+20\gamma H(3,1)+10z_2hh_2.
\end{align}
One derivation first writes the residue polynomial as
\begin{align*}
 r_{5,n}={}&-120e_4(n)+120\gamma e_3(n)
       +60(z_2-\gamma^2)e_2(n)+20A H_n+r_0\\
 &+20(H_n-\gamma)H_n^{(3)}
       +5\bigl(H_n^{(2)}\bigr)^2-25H_n^{(4)}+10z_2H_n^{(2)}.
\end{align*}
Apply \eqref{q5:finite-prefix} to every term, and then use the finite
product rule to collect the products occurring in \eqref{q5:S5finite}.

For the mixed residue there is a shorter identity.  Put
$C_0=4\gamma z_3+z_2^2-5z_4$.  Then
\begin{align}\label{q5:SIfinite}
 \sum_{n=1}^N\frac{r_{I,n}}n={}&
 4E_2(h_3-z_3)
 +h\{2H(2,2)-4h_4-4\gamma h_3+2z_2h_2+C_0\}\notag\\
 &-4H(3,1,1)-2H(2,1,2)-2H(2,2,1)+2H(2,3)\notag\\
 &+4H(4,1)+4h_5+4\gamma H(3,1)-4z_3h_2-2z_2H(2,1).
\end{align}
This is again a finite product-and-prefix calculation.

\begin{lemma}[All finite constants]
\label{q5:constants}
As $N\to\infty$, with $L=\log N$,
\begin{align}
 \sum_{n=1}^N\frac{r_{5,n}}n
   &=-L^5+20z_2L^3-90z_4L+C_5+o(1),
   \label{q5:S5asymptotic}\\
 r_{5,N}\log(N+1)
   &=-5L^5+60z_2L^3-90z_4L+o(1),
   \label{q5:r5asymptotic}
\end{align}
where
\begin{equation}\label{q5:C5}
 C_5=-\gamma^5+30\gamma^3z_2-80\gamma^2z_3
         -\frac{65}{2}\gamma z_4+60z_2z_3-69z_5.
\end{equation}
For the mixed residue,
\begin{equation}\label{q5:CI}
 \sum_{n\geq1}\frac{r_{I,n}}n
      =C_I=\gamma z_4-12z_2z_3+15z_5,
 \qquad r_{I,N}\log(N+1)=o(1).
\end{equation}
Finally the nonlogarithmic terms satisfy
\begin{equation}\label{q5:Bconstants}
 B_5-C_5=30\gamma^2z_2-60\gamma z_3+\frac{35}{3}z_4,
 \qquad
 B_I-C_I=-2\gamma z_3+\frac{35}{6}z_4.
\end{equation}
\end{lemma}
\begin{proof}
In \eqref{q5:S5finite}, substitute
$h=\log N+\gamma+O(N^{-1})$ and
$h_j=z_j+O(N^{1-j})$ for $j\geq2$, and replace the words with
outer exponent at least two by their limits in
\eqref{q5:MZVtable}.  Every discarded term is $o(1)$ even after
multiplication by the displayed powers of $\log N$.
The elementary polynomials are fixed, without any interpolation, by
\[
 E_0=1,\qquad
 jE_j=\sum_{r=1}^j(-1)^{r-1}h_rE_{j-r}.
\]
This gives \eqref{q5:S5asymptotic} including its constant
\eqref{q5:C5}.  For \eqref{q5:r5asymptotic}, directly substitute
$b_N=L+O(N^{-1})$, $c_N=2z_2+O(N^{-1})$,
$d_N=O(N^{-2})$, and $e_N=2z_4+O(N^{-3})$ in the residue.
Its constant polynomial term is
$-40z_2^2+10z_4=-90z_4$.

For \eqref{q5:SIfinite}, the first term tends to zero.  The expression
in braces multiplying $h$ has zero limit by
$Z_2(2,2)=3z_4/4$ and $z_2^2=5z_4/2$; its convergence is fast
enough that its product with $h$ tends to zero.  The remaining convergent
words have value
\begin{align*}
 &-4Z_3(3,1,1)-2Z_3(2,1,2)-2Z_3(2,2,1)+2Z_2(2,3)\\
 &\qquad+4Z_2(4,1)+4z_5+4\gamma Z_2(3,1)-6z_2z_3\\
 &\hspace{15mm}=\gamma z_4-12z_2z_3+15z_5.
\end{align*}
Also $r_{I,n}=O(n^{-1})+O(\log n/n^2)$, since its prospective
constant term $4z_2^2-10z_4$ vanishes.  This proves the ordinary
convergence and the boundary assertion in \eqref{q5:CI}.
Finally, expand the products in \eqref{q5:B5sum} and
\eqref{q5:BIsum} by the finite quasi-shuffle rule, prefix their
outer denominator, and use \eqref{q5:MZVtable}.
The results are respectively $C_5$ and $C_I$ plus the two explicitly
displayed corrections in \eqref{q5:Bconstants}.
\end{proof}

The preceding calculation is a full constant-term computation.  In
particular it retains the $-90z_4\log N$ terms on both sides of
the summation-by-parts boundary; neither term may be dropped before
their cancellation.

\subsection{Identification of the logarithmic terms}

The residue differences are particularly short.  For $n\geq1$,
with $b=b_{n-1}$, $c=c_{n-1}$, $d=d_{n-1}$,
\begin{align}
 r_{5,n}-r_{5,n-1}
    &=-\frac{20(b^3-3bc+d)}n+\frac{20b}{n^3}
                              +\frac{10c}{n^2},
      \label{q5:diff5}\\
 r_{I,n}-r_{I,n-1}
    &=\frac{4d}n+\frac{4b}{n^3}+\frac{2c}{n^2}.
      \label{q5:diffI}
\end{align}
For any residue sequence, finite summation by parts gives
\begin{equation}\label{q5:logSBP}
 \sum_{n=1}^Nr_{X,n}\log\left(1+\frac1n\right)
  =r_{X,N}\log(N+1)
       -\sum_{n=1}^N(r_{X,n}-r_{X,n-1})\log n.
\end{equation}
For $I$, equations \eqref{q5:integrated}, \eqref{q5:CI},
\eqref{q5:Bconstants}, and \eqref{q5:diffI} immediately give
\eqref{q5:Iarithmetic} and hence \eqref{q5:Ibridge}.

For $Q_5$, define
\[
 M_N=\sum_{n=1}^N\frac{(b_{n-1}^3-3b_{n-1}c_{n-1}+d_{n-1})\log n}{n}.
\]
Equations \eqref{q5:S5asymptotic}--\eqref{q5:logSBP} give
\begin{align}\label{q5:limit}
 Q_5={}&B_5-C_5+
       20\lim_{N\to\infty}
          \left(M_N-\frac{\log^5N}{5}+2z_2\log^3N\right)\notag\\
 &-20\sum_{n\geq1}\frac{b_{n-1}\log n}{n^3}
       -10\sum_{n\geq1}\frac{c_{n-1}\log n}{n^2}.
\end{align}
Splitting off $\log^3n-6z_2\log n$ and using the ordinary
Stieltjes limits gives \eqref{q5:Qarithmetic}.

For the exact spectral form introduce, initially for
$\operatorname{Re}s>1$,
\[
 T(s)=\sum_{n\geq1}
       \frac{b_{n-1}^3-3b_{n-1}c_{n-1}+d_{n-1}}{n^s}.
\]
Finite elementary-symmetric algebra gives
\begin{align}\label{q5:Tidentity}
 T(s)={}&6Z_4(s,1,1,1)-6\gamma Z_3(s,1,1)
           +3(\gamma^2-z_2)Z_2(s,1)\notag\\
 &-(\gamma^3-3\gamma z_2+2z_3)\zeta(s)-R(s).
\end{align}
Moreover $T(s)+\zeta'''(s)-6z_2\zeta'(s)$ is holomorphic in
$\operatorname{Re}s>0$.  Thus $T$ has principal part
$6u^{-4}-6z_2u^{-2}$ at $s=1+u$, and the finite limit in
\eqref{q5:limit} is exactly $-[u^1]T(1+u)$.
One can see the sign directly by differentiating the normally convergent
subtracted Dirichlet series; the ordinary limits contribute
$\gamma_4-6z_2\gamma_2$.
By \eqref{q5:Tidentity},
\[
 [u^1]T(1+u)=6\xi-6\gamma\theta+3(\gamma^2-z_2)\eta
       -(\gamma^3-3\gamma z_2+2z_3)\lambda_1-R'(1).
\]
The two remaining sums in \eqref{q5:limit} are
\[
 -\left.\frac{\partial}{\partial s}Z_2(s,1)\right|_{s=3}
          +\gamma\zeta'(3),\qquad
 -\left.\frac{\partial}{\partial s}Z_2(s,2)\right|_{s=2}
          -z_2\zeta'(2),
\]
respectively.  Substitution proves \eqref{q5:bridge}; subtraction
of five times \eqref{q5:Ibridge} gives \eqref{q5:QminusI}.

\subsection{A differential-polynomial interpretation and exact endpoints}

The fifth inverse-Gamma differential polynomial is
\begin{align}\label{q5:A5}
 \mathcal A_5
 ={}&\psi^5-10\psi^3\psi'+15\psi(\psi')^2+10\psi^2\psi''
            -10\psi'\psi''-5\psi\psi'''+\psi''''.
\end{align}
It equals $-\Gamma\,D^5(1/\Gamma)$.
Let $\Delta G=G(1)-[x^0]G(x)$ for a Laurent-meromorphic function
at zero.  Coordinate finite parts obey
$\operatorname{FP}\int_0^1G'(x)\,dx=\Delta G$.
The complete endpoint table needed for \eqref{q5:A5} is
\begin{equation}\label{q5:endpoints}
\begin{array}{c|c}
G&\Delta G\\\hline
\psi^4&12\gamma^2z_2-12\gamma z_3-11z_4\\
\psi^2\psi'&2\gamma z_3+3z_4/2\\
(\psi')^2&-6z_4\\
\psi\psi''&8z_4\\
\psi'''&0
\end{array}
\end{equation}
Integration by parts, with these constants retained, gives
\begin{align}\label{q5:Bell-explanation}
 Q_5-5I
   &=\operatorname{FP}\int_0^1\mathcal A_5(x)\,dx
       +30\gamma^2z_2-50\gamma z_3-\frac{35}{2}z_4.
\end{align}
Consequently \eqref{q5:QminusI} is compatible with the all-degree
inverse-Gamma construction, while the raw fifth power contains $I$.
This distinguishes a proved differential-polynomial cancellation from
the separate proposed cancellation by a scalar polynomial in $\psi$.
The formulas here neither prove nor disprove the existence of a further
ordinary-constant identity eliminating $J'(1)$.

\subsection{Normalized primitives and the polygamma hierarchy}

For $n\geq1$, define on $\operatorname{Re}z>0$
\[
 F_1(n,z)=\log\left(1+\frac zn\right)-\frac zn,
 \qquad
 F_j(n,z)=\frac{(n+z)^{1-j}-n^{1-j}}{1-j}-\frac z{n^j}
 \quad(j\geq2),
\]
with the logarithm real for positive arguments.

\begin{corollary}[Two normalized primitives]\label{q5:primitives}
For $X\in\{5,I\}$ the expression
\begin{align}\label{q5:primitiveformula}
 \mathcal P_X(z)={}&r_{X,0}\log z
       +\sum_{j=2}^5\frac{p_{X,j,0}z^{1-j}}{1-j}
       +\kappa_Xz
       +\sum_{n\geq1}\sum_{j=1}^5p_{X,j,n}F_j(n,z)
\end{align}
satisfies $\mathcal P_X'(z)=F_X(z)$ and has zero constant term
at the singular origin.  For real $b>0$,
\[
 \operatorname{FP}\int_0^bF_X(x)\,dx=\mathcal P_X(b).
\]
For $0<a<b$ the ordinary integral is
$\mathcal P_X(b)-\mathcal P_X(a)$.
\end{corollary}
\begin{proof}
Each $F_j$ is the primitive vanishing at zero of
$(n+z)^{-j}-n^{-j}$.  On compact sets it is $O(n^{-j-1})$;
the coefficient bounds from Lemma~\ref{q5:ML} therefore justify normal
convergence and differentiation.  The series converges uniformly on
each disk $|z|\leq r<1$, where all its summands vanish at zero.
The explicitly displayed singular terms and logarithm consequently
leave no constant term.
\end{proof}

\begin{corollary}[All fixed polygamma orders]\label{q5:derivative}
For every integer $r\geq1$, away from the poles,
\begin{equation}\label{q5:derivatives}
 D^rF_X(z)=(-1)^r r!
       \sum_{n\geq0}\sum_{j=1}^5
       \binom{r+j-1}{j-1}\frac{p_{X,j,n}}{(n+z)^{r+j}}.
\end{equation}
The series converges normally.  The left sides are explicitly
\[
 D^rF_5=
 \sum_{r_1+\cdots+r_5=r}\frac{r!}{r_1!\cdots r_5!}
                  \prod_{i=1}^5\psi^{(r_i)},
 \qquad
 D^rF_I=
 \sum_{r_1+r_2+r_3=r}\frac{r!}{r_1!r_2!r_3!}
         \psi^{(r_1)}\psi^{(r_2+1)}\psi^{(r_3+1)}.
\]
At $z=1$ these are finite polynomials in $\gamma$ and ordinary
zeta values, using
$\psi^{(j)}(1)=(-1)^{j+1}j!\zeta(j+1)$ for $j\geq1$.
\end{corollary}
\begin{proof}
Differentiate \eqref{q5:MLformula}.  For $r\geq1$ all constant
subtractions disappear, and the summands are bounded on compact sets
by a constant times $n^{-r-1}(1+\log^4n)$.  The product rule gives
the displayed left sides.
\end{proof}

\subsection{Verification and scope}

The accompanying exact replay checks 87 identities: principal
coefficients, constant Laurent terms, residue differences, both finite
harmonic summation formulas, all finite asymptotic constants,
nonlogarithmic reductions, spectral extraction, endpoint terms,
primitive derivatives, and the derivative kernels through order eight.
The finite harmonic checks are symbolic identities for arbitrary $N$;
they are not a fit to a finite sample of integers.  As regression checks,
the same summation code recovers the previously proved cubic and
quartic constants.

Independent numerical evaluations give
\begin{align*}
 Q_5&=-8.69030396262453826629558966236565357677029354492388429266\ldots,\\
 I&=-1.01946582194172926744369821391554613746742300351148475173\ldots.
\end{align*}
One route integrates the convergent Taylor series of the functions
after their fifth-order singularity has been factored out near zero,
and performs ordinary quadrature away from zero.  Cutoffs $1/8$ and
$1/16$ agree to about 77 decimal digits.  A second route uses
\eqref{q5:Iarithmetic} and \eqref{q5:Qarithmetic}, summing a finite
prefix and expressing a truncated digamma asymptotic tail through
Hurwitz-zeta derivatives.  At the finest recorded settings the
differences from the direct calculation are below
$1.02\cdot10^{-48}$ and $3.83\cdot10^{-48}$, respectively.
These are floating-point cross-checks, not interval-certified
enclosures.  The exact identities rest on the proofs above.

The advance relative to the inspected quartic report is the complete
fifth-power bridge, its mixed-coordinate identification, all finite
corrections, and the normalized primitive and derivative consequences.
The classical partial-fraction and multiple-zeta tools are not claimed
as new.  The scalar-polynomial elimination question remains distinct:
the presently proved cancellation naturally uses the additional
integral $I$ or, equivalently, the differential polynomial
$\mathcal A_5$.

\section{Polylogarithmic kernels, derivatives, and normalized primitives}
\label{sec:calculus}

The residue normalization has useful consequences beyond one finite
part. It gives convergent antiderivatives, strict multiple-zeta Taylor
coefficients, and explicit polylogarithmic kernels. These are exact
identities of functions on stated domains. The Gamma and polylogarithm
ingredients are classical; the point here is that the same normalization
works for the ordinary integral, the poles, and all subsequent
differentiations.

\subsection{A primitive with a fixed additive constant}

\begin{theorem}[The normalized Gamma-polygamma primitive]
\label{thm:Cprimitive}
Let $m\ge2$ and $\Omega_m=\C\setminus(-\infty,-(m-1)]$.
With the principal logarithm, put
\begin{equation}\label{eq:Cprimitive}
 \mathfrak P_m(z)=-m!\zeta(m)z
  -m!\sum_{n\ge m-1}e_{m-1}(n)
       \left\{\Log\left(1+\frac zn\right)-\frac zn\right\}.
\end{equation}
Then $\mathfrak P_m$ is holomorphic on $\Omega_m$,
$\mathfrak P_m(0)=0$, and $\mathfrak P_m'=\CM_m$. In particular
\begin{equation}\label{eq:primitive_atone}
 \mathfrak P_m(1)=m!(c_m-\lambda_{m-1}).
\end{equation}
For any real $-(m-1)<a<b$, the ordinary integral of $\CM_m$ is
$\mathfrak P_m(b)-\mathfrak P_m(a)$.
\end{theorem}
\begin{proof}
On every fixed compact set the bracket in \eqref{eq:Cprimitive} is
$O(n^{-2})$, uniformly in $z$. The logarithm is holomorphic on the
stated slit because the first nonzero residue occurs at $n=m-1$.
The normal convergence and the same bound for the differentiated
bracket justify termwise differentiation, giving \eqref{eq:CML}.
Every term vanishes at zero. Theorem~\ref{thm:Cintegral} gives the
value at one.
\end{proof}

\subsection{The all-ones polylogarithm as a kernel}

The harmonic generating identity is
\begin{equation}\label{eq:harmonic_polylog_kernel}
 \sum_{n\ge1}e_d(n)y^n
    =\frac{\Li_{\{1\}^d}(y)}{1-y}
    =\frac{(-\log(1-y))^d}{d!(1-y)},\qquad |y|<1.
\end{equation}
The factor $(1-y)^{-1}$ records the inclusive upper cutoff $n$.
Omitting it changes the harmonic convention and the resulting poles.

\begin{corollary}[Polylogarithmic integral and primitive kernels]
\label{cor:Cpolylog}
For $m\ge2$ and $\re z>-(m-1)$,
\begin{align}
 \CM_m(z)={}&-m!\zeta(m)
  +m!\int_0^1(1-y^z)
       \frac{\Li_{\{1\}^{m-1}}(y)}{y(1-y)}\,dy,
 \label{eq:Cpolylog}\\
 \mathfrak P_m(z)={}&-m!\zeta(m)z
  +m!\int_0^1\left\{z-\frac{y^z-1}{\log y}\right\}
       \frac{\Li_{\{1\}^{m-1}}(y)}{y(1-y)}\,dy.
 \label{eq:Ppolylog}
\end{align}
Here $y^z=e^{z\log y}$ uses the real logarithm for $0<y<1$.
Both integrals are ordinary absolutely convergent integrals.
\end{corollary}
\begin{proof}
Insert $(n+z)^{-1}-n^{-1}=\int_0^1y^{n-1}(y^z-1)\,dy$
into \eqref{eq:CML} and use \eqref{eq:harmonic_polylog_kernel}.
Near zero the polylogarithm is $O(y^{m-1})$, giving precisely
the half-plane in the statement. Near one, $1-y^z=O(1-y)$;
the remaining logarithmic power is integrable. These estimates are
uniform on compact parameter sets and justify summation and
differentiation. Integrating with respect to $z$ from zero gives
the second formula and its normalization.
\end{proof}

\subsection{All derivatives and the exact Taylor radius}

\begin{corollary}[Derivative Euler sums]\label{cor:Cderivatives}
For every $m\ge2$ and integer $r\ge1$,
\begin{equation}\label{eq:Cderivatives}
 \CM_m^{(r)}(z)=(-1)^{r+1}m!r!
       \sum_{n\ge m-1}\frac{e_{m-1}(n)}{(n+z)^{r+1}}.
\end{equation}
The series converges normally away from the actual poles. At zero,
\begin{equation}\label{eq:Cderivative_zero}
 \CM_m^{(r)}(0)=(-1)^{r+1}m!r!
  \left\{Z_m(r+1,\{1\}^{m-1})
       +Z_{m-1}(r+2,\{1\}^{m-2})\right\}.
\end{equation}
Equivalently, the Taylor series is
\begin{equation}\label{eq:Ctaylor}
 \CM_m(z)=-m!\zeta(m)
  +m!\sum_{r\ge1}(-1)^{r+1}
  \left\{Z_m(r+1,\{1\}^{m-1})
       +Z_{m-1}(r+2,\{1\}^{m-2})\right\}z^r.
\end{equation}
Its radius of convergence is exactly $m-1$.
\end{corollary}
\begin{proof}
Differentiate the normally convergent partial fractions. The bound
$e_{m-1}(n)=O((1+\log n)^{m-1})$ is sufficient for every
$r\ge1$. The exact finite identity
$e_d(n)=e_d(n-1)+e_{d-1}(n-1)/n$ separates the inclusive
cutoff into two strict sums, proving \eqref{eq:Cderivative_zero}.
The Taylor coefficients follow by division by $r!$. Finally the
origin is removable and the nearest actual singularity is the
simple pole at $-(m-1)$, whose residue is nonzero by
Theorem~\ref{thm:Cpoles}; this proves the exact radius.
\end{proof}

For example $\CM_2'(0)=4\zeta(3)$ and
$\CM_3'(0)=15\zeta(4)/2$. These recover the same coefficients
by differentiating the Bell-polygamma expressions. More generally,
every derivative at zero becomes a convergent multiple-zeta value,
whereas integration over one unit interval produces a regular
Laurent coefficient. This provides a direct functional connection
between the two types of constants.

On $\re z>-(m-1)$ the derivative also has the kernel
\begin{equation}\label{eq:Cderivative_kernel}
 \CM_m^{(r)}(z)=(-1)^{r+1}m!
  \int_0^1y^z(-\log y)^r
       \frac{\Li_{\{1\}^{m-1}}(y)}{y(1-y)}\,dy.
\end{equation}
The powers of $\log y$ therefore represent spatial differentiation,
while the logarithmic Mellin powers in \eqref{eq:alllaurent}
represent spectral differentiation.

The Gamma recurrence also supplies the exact discrete relation
\begin{equation}\label{eq:Cshift}
 \CM_m(x+1)-\CM_m(x)=\frac m x\CM_{m-1}(x+1),\qquad m\ge1.
\end{equation}
It follows by taking coefficients in
$G(x+1,t)=xG(x,t)/(x-t)$. The identity holds meromorphically;
removable values at exceptional points are obtained by continuation.
At positive integers it is the classical complete-harmonic recurrence.

\subsection{Antiderivatives in the spectral parameter}

The normalized germ $\Lambda_d$ from Corollary~\ref{cor:allcoefficients}
has a convergent primitive as well. On $\re u>-1$, define
\begin{equation}\label{eq:spectral_primitive}
 \mathfrak S_d(u)=\frac1{d!}\int_0^\infty\frac{V_d(v)}v
   \frac{(u\log v-1)v^u+1}{(\log v)^2}\,dv.
\end{equation}
At $v=1$ the quotient is assigned its limiting value $u^2/2$.
The same estimates as in Proposition~\ref{prop:mellin} justify
parameter differentiation and give
\begin{equation}\label{eq:spectral_primitive_derivative}
 \mathfrak S_d(0)=0,\qquad \mathfrak S_d'(u)=\Lambda_d(u).
\end{equation}
Indeed the quotient equals $\int_0^u wv^w\,dw$ on any path in the
half-plane. This identity integrates the completed spectral germ,
so its pole subtraction and integration constant are explicit.

For comparison, the ordinary Hurwitz Stieltjes constants are defined
by $\zeta(1+u,a)=u^{-1}+\sum(-1)^j\gamma_j(a)u^j/j!$,
with $\re a>0$. Their classical parameter calculus follows at once
from $\partial_a\zeta(s,a)=-s\zeta(s+1,a)$ \cite{DLMF}:
\begin{align}
 \partial_a\gamma_0(a)&=-\zeta(2,a),\\
 \partial_a\gamma_j(a)&=(-1)^{j+1}
     \{\partial_s^j\zeta(2,a)+j\partial_s^{j-1}\zeta(2,a)\},
       \qquad j\ge1,\label{eq:hurwitz_derivative}\\
 \int\gamma_j(a)\,da
     &=\frac{(-1)^{j+1}}{j+1}
                       \partial_s^{j+1}\zeta(0,a)+\text{constant}.
 \label{eq:hurwitz_primitive}
\end{align}
For the last identity expand
$\partial_a\zeta(u,a)=-u\zeta(1+u,a)$ at $u=0$.
These familiar formulas provide a sign and factorial control when
translating between ordinary Hurwitz constants and the strict-depth
constants used in the new finite-part identities.

% Standalone research fragment for the parent article; no preamble.
% Source: incoming Quartic_Tails_and_Complex_Resolvents, sections/mixed_even_tails.tex
% and sections/mixed_even_tail_generators.tex, source commit
% b7667ae40a2957eed9a53d8559d53a3ee20dd1c1.
\section{Finite reconstruction of every two-tail moment}
\label{tr:section}

The two-tail report proves a finite formula and a fixed coordinate space
for every centered moment, but asks whether a uniformly specified finite
prefix of the moments generates that space. We answer this question with
an explicit bound and a rational reconstruction matrix. For equal tail
orders the first $q$ moments suffice, with an exact determinant. For
unequal orders the bound below is deliberately not asserted to be optimal.
All ranks in this section are ranks of matrices over $\mathbb Q$; the
argument makes no assertion of arithmetic independence of zeta values.

The exact target is Question~8, ``Finite generation by initial moments'',
in \emph{Quartic Digamma Bridges, Harmonic Tail Moments, and Complex Resolvent Powers}. Its source theorem is
``All even orders and all centered moments'', with source label
\texttt{tail:all-moments}; the inverse transform has source label
\texttt{tail:inverse-transform}. The archive's Git blob is recorded
in the provenance manifest; it was inspected at the repository
commit printed on the title page.
The source formulas are reproduced below to fix every normalization;
the new argument begins with the finite-prefix theorem.

Let $2\le p\le q$ be even, $b=q/2$, and
\[
 A_p(n)=\zeta(p)-H_n^{(p)},\qquad
 D_r=D_{p,q}(-2r),\qquad
 D_{p,q}(s)=\sum_{n\ge0}A_p(n)A_q(n)(n+1/2)^{-s}.
\]
Outside its initial convergence half-plane, $D_{p,q}$ denotes its
meromorphic continuation. Every displayed $D_r$ is regular. The source
two-tail theorem gives blocks $C_j=C_{p,q}(2j+1)$ as follows. Put $W=p+q$
and use the decreasing-index convention
\[
 \zeta(a,b)=\sum_{n>k\ge1}n^{-a}k^{-b},\qquad a\ge2,\quad b\ge1.
\]
For a positive odd integer $d$, let
\begin{align}
 U_{p,q}(d)&=
 \begin{cases}
 \zeta(p,q-d)+\frac12\zeta(W-d),&d<q,\\[4pt]
 \displaystyle\sum_{h=1}^{\min\{(d-q+1)/2,p/2-1\}}
 \frac{\binom{d-q+1}{2h}}{d-q+1}B_{d-q+1-2h}\zeta(p-2h),&d>q,
 \end{cases}\label{tr:U}\\
 R_{p,q}(d)&=
 \begin{cases}
 0,&d<W-1,\\[3pt]
 \displaystyle-\frac{B_{d-W+1}}2
 \left\{\frac1p\binom{d-q}{p-1}+\frac1q\binom{d-p}{q-1}\right\},
 &d\ge W-1,
 \end{cases}\label{tr:R}\\
 C_j&=U_{p,q}(2j+1)+U_{q,p}(2j+1)+R_{p,q}(2j+1).
 \label{tr:Cdefinition}
\end{align}
Every sum is finite; the two alternatives in \eqref{tr:U} exhaust the
possibilities because $q$ is even. These blocks satisfy
\begin{align}
 D_r&=\sum_{j=0}^r T_{rj}C_j,
 &T_{rj}&=\frac{\binom{2r}{2j}}{2j+1}B_{2r-2j}(1/2),
 \label{tr:forward}\\
 C_j&=\sum_{r=0}^j B_{jr}D_r,
 &B_{jr}&=\binom{2j+1}{2r}2^{-2(j-r)}.
 \label{tr:backward}
\end{align}
The $B$ notation here denotes this finite matrix only when it has two
indices; $B_n(x)$ and $B_n$ are Bernoulli polynomials and numbers.

\paragraph{Why the inherited finite formula applies.}
For completeness, its finite-sum mechanism is short. The centered
Hurwitz expansion of $A_p(n)$ contains only the powers
$(n+1/2)^{1-p-2h}$. Hence $A_pA_q$ contains only even inverse powers;
asymptotic subtraction proves regularity of $D_{p,q}$ at $-2r$.
All removed nonnegative even powers have continued sum zero, since
$\zeta(-2h,1/2)=0$. Let $\operatorname{CT}_N$ mean the constant
coefficient at positive infinity and put
\[
 S_r(N)=\sum_{n=0}^{N-1}(n+1/2)^{2r}
       =\sum_{j=0}^r T_{rj}N^{2j+1}.
\]
This odd polynomial has no constant term. Finite summation by parts,
using $A_p(k-1)=A_p(k)+k^{-p}$, gives
\[
 D_r=\sum_{j=0}^rT_{rj}\operatorname{CT}_N E_{p,q,2j+1}(N),
\]
where
\[
 E_{p,q,d}(N)=N^dA_p(N)A_q(N)
 +\sum_{k=1}^N k^{d-q}A_p(k)
 +\sum_{k=1}^N k^{d-p}A_q(k)
 +\sum_{k=1}^N k^{d-W}.
\]
If $d<q$, the first ordered sum converges to $\zeta(p,q-d)$.
If $d>q$, write $m=d-q$ and
$F_m(N)=(B_{m+1}(N+1)-B_{m+1})/(m+1)$. The exact identity
\[
 \sum_{k=1}^Nk^mA_p(k)
 =A_p(N)F_m(N)+\sum_{\ell=1}^N\frac{F_m(\ell-1)}{\ell^p}
\]
reduces its constant term to the Bernoulli sum in \eqref{tr:U}.
The half-leading Faulhaber terms leave precisely one half of
$\zeta(W-d)$ for each convergent ordered block. At $d=W-1$ their
two $-\gamma/2$ constants cancel the final power sum's $+\gamma$.
The remaining three boundary products give exactly \eqref{tr:R},
by the uncentered Hurwitz expansion. This proves
$\operatorname{CT}_NE_{p,q,d}=C_{p,q}(d)$ and explains the fixed
finite-part normalization in \eqref{tr:forward}. Finally the ordinary
Bernoulli generating identity proves \eqref{tr:backward}. These are
inherited identities; no new priority is claimed for their derivation.

The source coordinate list is
\begin{equation}\label{tr:coordinates}
 C_0,\ldots,C_{b-1},\quad Z_0,\ldots,Z_{b-1},
 \qquad Z_k=\zeta(2k),\quad Z_0=-\frac12.
\end{equation}
It has $q$ specified entries. Using $Z_0$ instead of $1$ is an invertible
rational change of coordinates. For $j\ge b$ the source block formula
is equivalently
\begin{equation}\label{tr:A}
 C_j=\sum_{k=0}^{b-1}a_{jk}Z_k,
\end{equation}
where, with $d=2j+1$ and $e=d-p-q+2k+1$,
\begin{equation}\label{tr:Aentries}
 a_{jk}=
 \begin{cases}
 B_e\left\{
 \displaystyle\frac{\binom{d-q+1}{p-2k}}{d-q+1}
 +\displaystyle\frac{\binom{d-p+1}{q-2k}}{d-p+1}
 \right\},& e\ge0,\\[5pt]
 0,&e<0.
 \end{cases}
\end{equation}
In this formula a binomial term is omitted unless its lower index is at
least $2$ and no greater than its upper index. This qualification matters:
including a spurious lower index $0$ changes the formula. The rational
block is exactly the $k=0$ term, since $Z_0=-1/2$ and
$\binom{u}{v}/u=\binom{u-1}{v-1}/v$.

\begin{theorem}[A prescribed finite prefix]\label{tr:main}
If $p<q$, put
\begin{equation}\label{tr:N}
 N=\frac{q(q-p)}2+\frac q2+p-1.
\end{equation}
The matrix
\[
 M=(a_{jk})_{b\le j\le N,\ 0\le k<b}
\]
has column rank $b$ over $\mathbb Q$. Consequently
\begin{equation}\label{tr:span}
 \operatorname{span}_{\mathbb Q}\{D_0,D_1,\ldots\}
 =\operatorname{span}_{\mathbb Q}\{D_0,\ldots,D_N\}
 =\operatorname{span}_{\mathbb Q}
     \{C_0,\ldots,C_{b-1},Z_0,\ldots,Z_{b-1}\}.
\end{equation}
These are equalities of spans of actual numbers, regardless of any
linear dependence among those numbers. The coefficient system on the
$q$ formal labels in \eqref{tr:coordinates} has rank exactly $q$.
For $p=q$, all conclusions hold with $N=q-1$.
\end{theorem}

The pair-dependent bound is at most $q(q-1)/2+1$, so the first
$q(q-1)/2+2$ moments always suffice. The sharper diagonal assertion is
proved separately below. The existence theorem answers the source
finite-prefix question; it does not establish that the first $q$ moments
always suffice when $p<q$.

\subsection{A differential description of the coefficient columns}

Write $h(t)=\coth(t/2)$ and $m=q-p$. The source generator, or a direct
termwise beta integration of the Bernoulli series, gives column functions
\begin{equation}\label{tr:column}
 \begin{split}
 f_k(t)={}&
 \frac{1}{2(p-2k)!(q-1)!}
       \int_0^t u^{p-2k}(t-u)^{q-1}h(u)\,du\\
 &+\frac{1}{2(q-2k)!(p-1)!}
       \int_0^t u^{q-2k}(t-u)^{p-1}h(u)\,du.
 \end{split}
\end{equation}
The first summand is omitted if $p-2k<2$. Both integrands otherwise
have a removable endpoint singularity. Their Taylor series converge
for $|t|<2\pi$. The coefficient of $t^{2j+1}/(2j+1)!$ in $f_k$ is
$a_{jk}$ for every $j\ge b$; lower coefficients will not be used.
This follows from $h(u)=2\sum_{n\ge0}B_{2n}u^{2n-1}/(2n)!$ and
Euler's beta integral.

For a vector $v=(v_0,\ldots,v_{b-1})$ over $\mathbb C$, put
\[
 f_v=\sum_{k=0}^{b-1}v_kf_k,
 \qquad Q(t)=\sum_{k=0}^{b-1}v_k\frac{t^{q-2k}}{(q-2k)!}.
\]
Thus $Q$ is an arbitrary even polynomial of degree at most $q$ with
$Q(0)=0$. Differentiating the convolution integrals $q$ times gives
\begin{equation}\label{tr:operator}
 2f_v^{(q)}(t)
 =\left(Q^{(m)}(t)-Q^{(m)}(0)\right)h(t)
       +(Q(t)h(t))^{(m)}
 =:\mathcal L_Q(t).
\end{equation}
The subtraction $Q^{(m)}(0)$ enforces the omitted zero lower index in
\eqref{tr:Aentries}. Each side is an odd function holomorphic at zero.

\begin{lemma}[A finite jet forces the polynomial to vanish]
\label{tr:zerobound}
Suppose $m=q-p>0$ and
\[
 \operatorname{ord}_{t=0}\mathcal L_Q\ge2N+3-q
\]
with $N$ as in \eqref{tr:N}. Then $Q=0$.
\end{lemma}

\begin{proof}
Set $y=e^t$. For $r\ge1$, direct differentiation shows that
$h^{(r)}(t)$ has the form $P_r(y)/(y-1)^{r+1}$, where $P_r$ is a
polynomial of degree at most $r$, divisible by $y$. Indeed it is true
for $h'=-2y/(y-1)^2$, and application of $y\,d/dy$ preserves the
asserted form. Also $h=(y+1)/(y-1)$.

It follows by the product rule that
\[
 E(t)=(e^t-1)^{m+1}\mathcal L_Q(t)
     =\sum_{\nu=0}^{m+1}P_\nu^*(t)e^{\nu t}
\]
is an entire exponential polynomial. For $\nu=0$ and $\nu=m+1$,
$\deg P_\nu^*\le p$; for $1\le\nu\le m$, the degrees are at most
$q$. To verify the endpoint degree statement, all terms involving
$h^{(r)}$ with $r\ge1$ have an $e^t$ factor and have exponential
degree at most $m$, after the denominator is cleared. Only the
$h$ term reaches exponent $0$ or $m+1$, and its polynomial factor
is $2Q^{(m)}-Q^{(m)}(0)$, of degree at most $p$.

Therefore $E$ is annihilated by the monic constant-coefficient
differential operator
\begin{equation}\label{tr:ode}
 \mathcal P(\partial_t)=
 \partial_t^{p+1}
 \left(\prod_{\nu=1}^{m}(\partial_t-\nu)^{q+1}\right)
 (\partial_t-(m+1))^{p+1}.
\end{equation}
Its order is
$H=2(p+1)+m(q+1)$. A holomorphic solution of this equation whose
first $H$ Taylor coefficients vanish is identically zero: the monic
equation recursively expresses each higher derivative in the previous
$H$ derivatives. On the other hand, the hypothesis gives
\[
 \operatorname{ord}_0E\ge(2N+3-q)+(m+1)
    =2N+4-p=H.
\]
Hence $E=0$ and $\mathcal L_Q=0$ meromorphically.

At every nonzero $\tau=2\pi i\ell$, $h(t)$ has residue $2$.
Since the even integer $m$ is at least $2$, the pole of order $m+1$
in $(Qh)^{(m)}$ has coefficient $2m!Q(\tau)$. The other term of
$\mathcal L_Q$ has at most a simple pole. Thus $Q(\tau)=0$ for all
nonzero integers $\ell$. A polynomial with infinitely many zeros is
zero. This proves the lemma.
\end{proof}

\begin{proof}[Proof of Theorem~\ref{tr:main} for unequal orders]
If $Mv=0$, the odd Taylor coefficients of $f_v$ at degrees
$q+1,q+3,\ldots,2N+1$ vanish. Differentiation $q$ times kills every
lower coefficient; the next possible degree is $2N+3-q$.
Equation~\eqref{tr:operator} and Lemma~\ref{tr:zerobound} give $Q=0$,
hence $v=0$. The matrix has rank $b$ already over $\mathbb C$, and
therefore over $\mathbb Q$.

Its finite left inverse recovers all $Z_k$ from $C_b,\ldots,C_N$.
Equations~\eqref{tr:forward}--\eqref{tr:backward} recover these blocks
and the initial $C_0,\ldots,C_{b-1}$ from $D_0,\ldots,D_N$.
All later blocks are given by \eqref{tr:A}; this proves all the span
equalities. The early $b$ coefficient rows are the identity on the
$C_0,\ldots,C_{b-1}$ labels, and the later rows have rank $b$ on
the remaining labels. Thus the full formal rank is $2b=q$.
\end{proof}

\subsection{An explicit rational reconstruction matrix}

The preceding proof is effective without a search for a nonsingular
minor. Define
\begin{equation}\label{tr:leftinverse}
 L=(M^{\mathsf T}M)^{-1}M^{\mathsf T}
   =\frac{\operatorname{adj}(M^{\mathsf T}M)M^{\mathsf T}}
          {\det(M^{\mathsf T}M)}.
\end{equation}
Every entry is an explicitly specified rational number obtained from
the finite Bernoulli formula. The denominator is nonzero because $M$
has full column rank; equivalently its real Gram matrix is positive
definite. If $V=(C_b,\ldots,C_N)^{\mathsf T}$, then
\begin{equation}\label{tr:reconstruct}
 (Z_0,\ldots,Z_{b-1})^{\mathsf T}=LV.
\end{equation}
Together with \eqref{tr:backward}, this is an explicit rational formula
in the required prefix of actual moments. For any desired later $D_r$,
insert \eqref{tr:reconstruct} into \eqref{tr:A} and then
\eqref{tr:forward}. This is a terminating exact algorithm, including
when the actual coordinate values satisfy unknown rational relations.

One can also extract exactly $q$ specified moments from this bounded
prefix. Form their formal coefficient rows $W_r\in\mathbb Q^q$ by
\[
 (W_r)_j=T_{rj}\quad(0\le j<b),\qquad
 (W_r)_{b+k}=\sum_{j=b}^r T_{rj}a_{jk}\quad(0\le k<b),
\]
where a term with $j>r$ is zero. Scan $r=0,1,\ldots,N$, retaining a
row exactly when it increases rational rank. The theorem proves that
this deterministic finite procedure retains $q$ rows, with indices
$r_1<\cdots<r_q\le N$. If $A_*$ is the matrix of those rows, then
\[
 D_r=W_rA_*^{-1}(D_{r_1},\ldots,D_{r_q})^{\mathsf T}
\]
for every $r\ge0$. This is an explicit basis-selection algorithm for
the formal system, not a claim that its $q$ actual numbers form an
arithmetically independent set.

\subsection{Equal orders: an exact determinant and a shorter prefix}

If $p=q=2b$, take the square late matrix with $b\le j\le2b-1$.
Write $j=b+t$ with $0\le t<b$. Formula~\eqref{tr:Aentries} is zero
whenever $t+k<b-1$, and on the antidiagonal $t+k=b-1$ it equals
$1/(t+1)$. Therefore
\begin{equation}\label{tr:diagonal-det}
 \det M=\frac{(-1)^{b(b-1)/2}}{b!}\ne0.
\end{equation}
This proves the asserted $N=q-1$. It also gives a particularly simple
reconstruction, for $t=0,\ldots,b-1$:
\begin{equation}\label{tr:diagonal-recursion}
 Z_{b-1-t}=(t+1)\left(
 C_{b+t}-\sum_{k=b-t}^{b-1}a_{b+t,k}Z_k\right).
\end{equation}
The sum is empty at $t=0$. In the coordinate ordering
\eqref{tr:coordinates}, the determinant of the complete first-$q$
moment matrix is
\begin{equation}\label{tr:complete-det}
 \frac{(-1)^{b(b-1)/2}}{b!\,(2q-1)!!}.
\end{equation}
Here the extra denominator is the product of the diagonal entries
$1/(2r+1)$ of the forward Bernoulli transform.

\subsection{A new concrete mixed example: orders two and six}

For $(p,q)=(2,6)$, the general theorem supplies a longer guaranteed
prefix, but direct exact elimination proves that the first six moments
already suffice. The last three block rows are
\[
 \begin{pmatrix}C_3\\C_4\\C_5\end{pmatrix}
 =\begin{pmatrix}
  2/3&5/12&-1/12\\
  5/6&-7/24&1/12\\
  -47/60&1/2&-3/20
 \end{pmatrix}
 \begin{pmatrix}-1/2\\\zeta(2)\\\zeta(4)\end{pmatrix},
\]
whose determinant is $61/5760$. Its inverse is
\begin{equation}\label{tr:26inverse}
 \frac1{61}
 \begin{pmatrix}
 12&120&60\\344&-952&-720\\1084&-3800&-3120
 \end{pmatrix}.
\end{equation}
For example, the rational row becomes the exact identity
\begin{equation}\label{tr:26rational}
 61D_0+8204D_1+116000D_2+370048D_3
      +303360D_4+56320D_5=-\frac{7808}{3}.
\end{equation}
It has the ordinary absolutely convergent interpretation
\begin{equation}\label{tr:26ordinary}
 \begin{split}
 \sum_{n\ge0}\bigg\{&P(x_n)A_2(n)A_6(n)
   -11264x_n^4-\frac{136960}{3}x_n^2-\frac{88064}{5}\bigg\}
   =-\frac{7808}{3},\\
 P(x)={}&61+8204x^2+116000x^4+370048x^6
                +303360x^8+56320x^{10},\qquad x_n=n+\frac12.
 \end{split}
\end{equation}
Indeed $A_2(n)A_6(n)=\frac15x_n^{-6}
-\frac4{15}x_n^{-8}+\frac{87}{200}x_n^{-10}+O(x_n^{-12})$,
so the displayed bracket is $O(n^{-2})$. Also
$A_2(n)A_6(n)=\psi'(n+1)\psi^{(5)}(n+1)/120$, giving a
polygamma identity with exactly the same normalization.

\paragraph{What is still open.}
The first-$q$ square coefficient matrix is nonsingular in every
unequal pair checked by exact rational arithmetic through $q=24$.
This is finite evidence, not a proof for all $p,q$. The general
theorem gives an explicit finite prefix and an exact rational inverse
without assuming that stronger assertion. It also leaves untouched
the arithmetic dimension of the actual number space: formal rank $q$
is compatible with unknown rational relations among its zeta coordinates.

\section{Prior art, coefficient conventions, and the source audit}
\label{sec:normalizer-prior-art-audit}

\subsection{Established ingredients and the attribution boundary}

The normalized Gamma generator is closely connected with the
classical Roman harmonic numbers.  In the notation of
Sesma~\cite[Proposition~3(a), equation~(49)]{Sesma},
\begin{equation}
 \frac{n!}{(1-t)_n}=\sum_{m\geq0}c_n^{(m)}t^m
       =\frac{\Gamma(1-t)\Gamma(n+1)}{\Gamma(n+1-t)}.
 \label{eq:audit-roman-generator}
\end{equation}
Consequently $\mathcal C_m(n+1)/m!=c_n^{(m)}$ is the complete harmonic
star sum with $m$ unit indices.  The complete harmonic-polynomial
formula and the beta-logarithmic integral are also recorded in
\cite[equations~(29)--(30)]{Sesma}; the relation of negative Roman
degrees to strict harmonic sums appears in equations~(24) and~(26).
The meromorphic functions $\mathcal C_m(x)$ used here are therefore natural
continuations of established harmonic Bell polynomials.  Their
generator and their positive-integer specialization should not
be presented as newly discovered objects.

The first ordinary-integral case is likewise known.
Coffey~\cite[Proposition~3(b), equation~(1.19)]{Coffey} gives
\begin{equation}
 \gamma_1=\zeta(2)+\int_0^1
       \left(\gamma\psi(x)
          +\frac{\psi(x)^2-\psi'(x)}{2}\right)\,dx.
 \label{eq:audit-coffey-first-case}
\end{equation}
Since $\mathcal C_2=(\psi+\gamma)^2+\zeta(2)-\psi'$, this is exactly
\begin{equation}
 \int_0^1 \mathcal C_2(x)\,dx
       =\gamma^2-\zeta(2)+2\gamma_1
       =2(c_2-\lambda_1),
 \label{eq:audit-coffey-normalized-case}
\end{equation}
where $c_k=[u^k]\Gamma(1+u)^{-1}$ and
$\lambda_1=[u]\zeta(1+u)=-\gamma_1$.
Coffey's Appendix~B also develops complex-degree Gamma derivatives
and their Bell-polynomial and polygamma descriptions.

Young's work on height-one multiple zeta functions establishes
the reciprocal-Gamma singular Laurent data at $s=1$ and
expressions for the first regular coefficient in terms of
Bernoulli numbers of the second kind and multiple Hurwitz
integrals~\cite{Young2023,Young2024,YoungLecture}.  It also relates
strict-depth Laurent constants to Ramanujan sums of complete
harmonic star sums.  In particular, the Cauchy-number series
for the linear coefficient in this article is existing
literature in equivalent notation.  To make the equivalence
explicit, put
\begin{equation}
 \lambda_d=[u]\zeta(1+u,\{1\}^{d-1}),\qquad
 a_n=\frac{1}{(n-1)!}\int_0^1(r)_n\,dr.
 \label{eq:audit-lambda-and-cauchy}
\end{equation}
If $t/\log(1+t)=\sum_{n\geq0}b_nt^n$, then
\begin{equation}
 \frac{a_n}{n}=\sum_{k>n}|b_k|,\qquad
 \sum_{n\geq1}\frac{a_n}{n^{d+2}}
      =\sum_{k\geq1}|b_k|H_{k-1}^{(d+1)}.
 \label{eq:audit-cauchy-equivalence}
\end{equation}
The first equality follows by comparing coefficients in
\begin{equation}
 \sum_{n\geq0}\frac{z^n}{n!}\int_0^1(r)_n\,dr
   =\frac{-z}{(1-z)\log(1-z)}
   =\frac{1-\sum_{k\geq1}|b_k|z^k}{1-z}.
 \label{eq:audit-cauchy-generating-equivalence}
\end{equation}
Here $\sum_{k\geq1}|b_k|=1$, obtained by letting $z\uparrow1$
in $-z/\log(1-z)=1-\sum_{k\geq1}|b_k|z^k$.
The second equality in \eqref{eq:audit-cauchy-equivalence}
then follows by Tonelli's theorem.  Thus
\begin{equation}
 \lambda_d=c_{d+1}+\sum_{n\geq1}\frac{a_n}{n^{d+2}}
 \label{eq:audit-known-cauchy-coefficient}
\end{equation}
is precisely Young's coefficient formula after a change of
Cauchy-number convention, rather than an independent new evaluation.

The contribution developed here is the explicit residue organization
needed for the finite-part problem: the scalar Gamma--Appell
normalizer, its uniqueness up to constants, its finite decaying
defect, and its comparison with the exact differential normalizer.
The ordinary integral represents the established strict-depth
coefficient in terms of Gamma and polygamma functions.
These statements are proved directly in the article.  The source
search supplies attribution for their ingredients; it does not
establish worldwide priority for every possible reformulation.

\subsection{A sign description in Young's 2024 paper}

There is a specific convention issue worth correcting before the
formulas are imported into a symbolic library.  In the standard
Stieltjes convention,
\begin{equation}
 \zeta(1+u)=\frac1u+\gamma-\gamma_1u+O(u^2).
 \label{eq:audit-standard-stieltjes-sign}
\end{equation}
Thus the literal linear coefficient is $-\gamma_1$.
Young's earlier lecture explicitly defines the regular coefficients
with the factors $(-1)^i/i!$ and, on its numbered slides~3 and~6,
correctly identifies $-\gamma_1^{[d]}(1)$ as the linear
coefficient~\cite{YoungLecture}.

In \cite[equations~(4.1)--(4.2) and Corollary~3]{Young2024},
the verbal description identifies the same signed Stieltjes
constant with the literal linear coefficient.  With the explicit
notation \eqref{eq:audit-lambda-and-cauchy}, the unambiguous
translation of equation~(4.2) is
\begin{equation}
 \nu_{0,d}-\nu_{-1,d-1}=c_{d+1}-\lambda_d,
 \label{eq:audit-young-sign-translation}
\end{equation}
where the two series in the source's notation are
\begin{equation}
 \nu_{0,d}=\sum_{n\geq2}\frac{(-1)^n}{n}
                       \zeta(n,\{1\}^{d}),\qquad
 \nu_{-1,d-1}=\sum_{n\geq2}\frac{(-1)^n}{n-1}
                       \zeta(n,\{1\}^{d-1}).
 \label{eq:audit-young-nu-definition}
\end{equation}
The second series is interpreted as its ordinary convergent
alternating sum when $d=1$; the other cases here are absolutely
convergent.  At depth one,
\begin{equation}
 \nu_{0,1}-\nu_{-1,0}-c_2=\gamma_1=-\lambda_1,
 \label{eq:audit-depth-one-sign-control}
\end{equation}
which exhibits the sign directly.  The appropriate correction is
to describe $\gamma_1^{[d]}(1)$ as the \emph{negative} of the
linear coefficient, or to translate every formula consistently
into the literal-coefficient notation $\lambda_d$.  The equations
for the conventional signed constants need not be reversed.

The finite-part convention is a separate issue.  In the reviewed
quartic-tail and Barnes incoming articles, the quartic constants
denoted $\theta$ and $\beta_1$ represent the same strict coefficient
$[u]\zeta(1+u,1,1)$, and the finite parts use a cutoff in $x$
itself.  No additional incompatibility between those conventions
was established in the audit.  A change of cutoff coordinate,
however, must be transported explicitly before two finite-part
formulas are compared; a change of Laurent-coefficient sign must
not be used to absorb such a coordinate change.

\section{Further research questions and concrete next targets}
\label{sec:questions}

The following programme emphasizes exact identities and structural
reduction. Its first two items sharpen source questions partially
answered here; the subsequent items follow from the proved formulas.
Numerical experiments may guide these questions, but their desired
outputs are identities, rank theorems, and analytic proofs.

\subsection{The first \texorpdfstring{$q$}{q} moments for unequal even orders}

\begin{conjecture}[The optimal consecutive prefix]\label{conj:firstq}
For all even $2\le p<q$, the square matrix
\[
 (a_{jk})_{q/2\le j\le q-1,\ 0\le k<q/2}
\]
is nonsingular over $\Q$. Equivalently, the first $q$ centered
moments have the full formal rank $q$.
\end{conjecture}

The article proves a larger finite prefix for every pair and the
conjectured prefix on the diagonal. Exact calculations verify the
unequal assertion through $q=24$, but this is a finite test.
A useful next proof would replace the general exponential-polynomial
zero count by a sharper property of the restricted even-polynomial
space in \eqref{tr:operator}. An explicit determinant formula would
also explain the nontrivial rational numerators already visible for
$(2,4)$ and $(2,6)$. A sign pattern alone would be insufficient unless
it proves nonvanishing for every order.

\subsection{Can the mixed fifth coordinate be eliminated?}

Theorem~\ref{q5:main} singles out the holomorphic quantity $J'(1)$.
The first unresolved question is whether it has an exact reduction
to the height-one coefficients and ordinary zeta/Stieltjes constants
already present in that theorem. It is not appropriate to declare a
new independent constant merely because it appears in one reduction.
Conversely, a small numerical integer relation would not establish
its reducibility.

A more algebraic version asks for a scalar polynomial in $\psi$
whose coordinate finite part contains only a prescribed highest-depth
coefficient and ordinary constants. The scalar Gamma--Appell theorem
controls the nondecaying residue polynomial at infinity. Its explicit
defects, starting at degree four, show why this statement alone does
not settle all convergent lower-depth contributions. A promising
next target is to classify the residue defects modulo discrete
differences whose logarithmic sums reduce by finite harmonic algebra.

\subsection{An all-degree classification of the mixed coordinates}

The differential sequence $\CM_m$ has exactly one height-one residue
word. The raw scalar power $\psi^m$ introduces products involving
higher polygammas. At degree five the additional coordinate is represented
by $\FP\int\psi(\psi')^2$. At degree six and beyond, one can ask:
which differential monomials survive integration by parts after every
finite endpoint constant is retained?

A precise project is to grade monomials by
$\operatorname{wt}(\psi^{(r)})=r+1$, quotient them by total
derivatives, and compute the finite-dimensional residue-coordinate
image at each weight. This must keep endpoint evaluation as part of
the map: discarding a total derivative without its finite constant
would produce a different problem. The desired result is a recursion
for the surviving spectral words, with a proved coordinate count,
rather than a list of fitted numerical bases.

\subsection{Shifted and rational-cutoff Gamma normalizations}

Replacing the base point $1$ by a parameter $a$ naturally produces
$\Gamma(a-t)\Gamma(x)/(\Gamma(a)\Gamma(x-t))$ and differences
of Hurwitz polygammas. The special base point in this article makes
the origin residue exactly linear in $t$, which is responsible for
the ordinary integrals at every degree $m\ge2$. For general $a$, a
new subtraction is required.

Find the canonical joint subtraction in $x$ and $t$, prove the
resulting integral bridge, and express it in the ordered
multiple-Hurwitz Laurent coordinates already developed in the
incoming Barnes report. Rational $a$ should connect the resulting
formulas with congruence-restricted harmonic sums and cyclotomic
polylogarithms. The question is the exact compatibility of the
normalizations, not just formal substitution in an unshifted formula.

\subsection{Four tails and higher products}

For products $\prod_{i=1}^k(\zeta(p_i)-H_n^{(p_i)})$ with even
$p_i$ and even $k$, the centered even-power asymptotics persist.
Finite summation by
parts produces higher-depth convergent zeta coordinates. Determine a
finite formal coordinate system and prove an effective initial-moment
generation theorem for even $k\ge4$. Odd numbers of tails instead
produce odd centered powers and can have poles at the target spectral
values; their finite-part moments constitute a separate extension.
The two-tail proof suggests a route
through differential operators on products of $\coth(t/2)$, but
higher pole multiplicities and overlaps of coordinate columns must
be handled explicitly. Merely observing stabilization in low-order
matrices would not prove a uniform theorem.

\subsection{Exact primitives at algebraic or rational arguments}

The functions $\mathfrak P_m$ provide explicit normally convergent
logarithmic primitives on a slit plane. Evaluate their values at
rational arguments and their reflection or multiplication combinations.
One may seek reductions to Barnes multiple-Gamma functions, multiple
Hurwitz derivatives, or cyclotomic polylogarithms. The removable values
\eqref{eq:Cnegative} and the shift identity \eqref{eq:Cshift} supply
exact boundary conditions. Any proposed primitive identity should
match both its derivative and one normalization value; matching
derivatives alone leaves an undetermined constant.

\subsection{Higher Laurent jets and their functional relations}

Equation~\eqref{eq:alllaurent} supplies every regular Laurent
coefficient through a fixed logarithmic Mellin kernel. Determine
which identities among those logarithmic moments arise from
polylogarithmic functional equations, shuffle products, and the
spectral primitive \eqref{eq:spectral_primitive}. A particularly
concrete question is whether selected mixed derivatives in the
sixth-power calculation can be expressed using only finitely many
height-one jets plus a small collection of higher-depth germs.

\subsection{Proof certificates and unresolved ladder identities}

The finite harmonic identities in the fifth-power ledger are suitable
for independent formal verification: only exact quasi-shuffle
multiplication, the prefix rule, explicit low-weight reductions, and
polynomial constant-term extraction are needed. The analytic layer
can be separated into normal convergence, the contour remainder,
and the Mellin continuation lemmas. This provides a manageable route
to integrating proof certificates into ProveIt.

The manuscript's remaining $S_6$ and revised $S_8$ ladder candidates
continue to deserve exact functional-equation research. Their
existing high-precision certificates should guide the search for
an identity or an obstruction, while remaining explicitly distinct
from a proof of equality. No result of this article changes their
status.

\section{Verification, provenance, and integration}
\label{sec:reproducibility}

The accompanying archive is an editable research package. It contains
this article in modular and single-file \LaTeX, the compiled PDF,
the exact calculation scripts, numerical evaluators, their recorded
JSON outputs, a build script, a checksum manifest, and proposed
repository integration notes. The scripts require Python, SymPy, and
mpmath; no network access or private data is needed to replay them.

\subsection{What the calculations verify}

The fifth-power exact replay checks 87 identities using finite harmonic
quasi-shuffle algebra. They include all principal coefficients, the
constant Laurent coefficients at zero, both finite residue-sum
identities, their asymptotic constants, the nonlogarithmic sums,
the spectral substitution, and every endpoint entry used in the
integration by parts. These are symbolic finite identities for
arbitrary cutoff $N$, together with explicitly stated classical
multiple-zeta reductions; the computation does not infer the
identities by sampling $N$.

The tail reconstruction replay checks 1,582 exact assertions. These
include the forward and inverse Bernoulli transforms, the agreement
of the compact coefficient matrix with the inherited finite formula,
the diagonal determinant, the displayed rational inverse, the
orders-$(2,6)$ identity, and finite instances of the
exponential-polynomial degree counts. It also verifies the first-$q$
square matrix for all 78 even pairs $2\le p\le q\le24$. Those
finite checks support the stronger conjecture but are not its proof.
The all-pairs finite-prefix theorem rests on the differential and
pole argument in Section~\ref{tr:section}.

The independent Gamma-normalizer replay checks 202 exact assertions
through degree eight and integer pole labels zero through eight,
including the removable values and the first nonzero residues.
Separate ordinary-integral computations for $m=2,\ldots,7$ use
analytic Gamma coefficients on $[0,1]$ and the completed logarithmic
Mellin kernel on $(0,\infty)$. The largest recorded absolute
discrepancy is below $7.403\cdot10^{-58}$ at 70 working decimal
digits. A local Taylor switch is changed independently to test
cancellation handling. The scalar normalization has its own 94-check finite polynomial
replay, including the inverse-Gamma finite-part convolution through
degree ten. Together the four exact suites verify 1,965 assertions.

\begin{center}
\small
\begin{tabular}{@{}>{\raggedright\arraybackslash}p{.29\textwidth}>{\raggedright\arraybackslash}p{.29\textwidth}>{\raggedright\arraybackslash}p{.31\textwidth}@{}}
\toprule
Independent comparison & Finest recorded absolute discrepancy & Type \\
\midrule
Fifth-power finite part & $<3.83\cdot10^{-48}$ & Taylor-subtracted
quadrature versus harmonic arithmetic with asymptotic tail \\[4pt]
Mixed integral $I$ & $<1.02\cdot10^{-48}$ & The same two independent
analytic representations \\[4pt]
Orders-$(2,6)$ rational sum & $<2.40\cdot10^{-62}$ & Hurwitz-tail
sum with separately derived centered tail \\[4pt]
Gamma ordinary integrals & $<7.403\cdot10^{-58}$ & Bell/Gamma
coefficients versus completed Mellin kernel \\
\bottomrule
\end{tabular}
\end{center}

All entries in this table are floating-point diagnostics. They are
not interval certificates and are not presented as proofs. In
particular, the Gamma evaluator records a bound for its omitted
analytic kernel tail, but that bound does not enclose quadrature,
local Taylor truncation, or rounding error. The mathematical
conclusions follow from the analytic and exact algebraic arguments
in the text.

\subsection{Source and editorial controls}

The provenance manifest records the inspected repository commit,
the source archive and relevant member hashes, and the external
primary references. The main manuscript was checked for prior
results before the new targets were selected. The audit does not
purport to certify every page of the repository. In particular,
the already proved Choi harmonic conjectures and the proved $S_4$
identity are not rediscovered here, and $S_6$ and the revised $S_8$
remain explicitly open.

The integration notes propose where to place the new sections and
how to update the three source questions. They also provide the
precise Stieltjes sign translation, so an importer can retain the
literal-coefficient notation consistently. No repository file has
been silently altered, and the package is intended for review and
integration through the project's normal process.

The broader literature also contains harmonic-zeta Stieltjes
families with different indexing conventions; see, for example,
Karg\i n--Dil--Cenkci--Can \cite{Kargin}. All uses of such constants
in future extensions should name the defining germ and its
factorial and sign convention explicitly, as done here.

\begin{thebibliography}{99}
\raggedright

\bibitem{Repo}
V. Reshetnikov, \emph{ProveIt}, GitHub repository, manuscript and
incoming reports inspected at commit
\texttt{b7667ae40a2957eed9a53d8559d53a3ee20dd1c1},
11 October 2026.
\url{https://github.com/VladimirReshetnikov/ProveIt}.
Repository paths and file hashes are recorded in the accompanying
\texttt{provenance/} directory.

\bibitem{Quartic}
\emph{Quartic Digamma Bridges, Harmonic Tail Moments, and Complex Resolvent Powers}, ProveIt incoming report,
11 October 2026. Archive
\texttt{ProveIt\_Quartic\_Tails\_Resolvents\_2026-10-11.zip},
in \texttt{docs/incoming/} at the snapshot in \cite{Repo}.
Principal source sections: \texttt{mixed\_even\_tails.tex},
\texttt{mixed\_even\_tail\_generators.tex}, and
\texttt{research\_questions.tex}.

\bibitem{Barnes}
\emph{Barnes Resolvents and Higher Harmonic Identities}, ProveIt incoming
report, 11 October 2026, in the incoming directory at the snapshot
in \cite{Repo}. Used for the height-one Laurent normalization and
the comparison of quartic finite-part conventions.

\bibitem{Young2023}
P. T. Young, Global series for height 1 multiple zeta functions,
\emph{European Journal of Mathematics} \textbf{9} (2023), article 99.
\url{https://doi.org/10.1007/s40879-023-00695-0}.
The publisher's abstract and metadata were inspected; the full
subscription article was not obtained. The public lecture
\cite{YoungLecture} was read in full.

\bibitem{Young2024}
P. T. Young, Series of height one multiple zeta functions,
\emph{Integers} \textbf{24} (2024), A43.
\url{https://doi.org/10.5281/zenodo.11221606}.
Public full text: \url{https://math.colgate.edu/~integers/y43/y43.pdf}.

\bibitem{YoungLecture}
P. T. Young, \emph{Global series for height 1 multiple zeta functions},
public lecture slides, 4 July 2023, 17 numbered slides with overlays.
\url{https://iecl.univ-lorraine.fr/wp-content/uploads/2023/07/tuesday-12-young-p.pdf}.
The Laurent convention occurs on numbered slides 3 and 6; complete
harmonic star sums and their Ramanujan sums occur on slides 16 and 17.

\bibitem{Sesma}
J. Sesma, The Roman harmonic numbers revisited,
\emph{Journal of Number Theory} \textbf{180} (2017), 544--565.
\url{https://doi.org/10.1016/j.jnt.2017.05.009}.
Public preprint: \url{https://arxiv.org/abs/1702.03718}.

\bibitem{Coffey}
M. W. Coffey, Series representations of the Riemann and Hurwitz zeta
functions and series and integral representations of the first
Stieltjes constant, arXiv:1106.5146, 2011.
\url{https://arxiv.org/abs/1106.5146}.

\bibitem{Kargin}
L. Karg\i n, A. Dil, M. Cenkci, and M. Can,
On the Stieltjes constants with respect to harmonic zeta functions,
\emph{Journal of Mathematical Analysis and Applications} (2023),
article 127302.
\url{https://doi.org/10.1016/j.jmaa.2023.127302}.
Public preprint: \url{https://arxiv.org/abs/2304.03517}.

\bibitem{DLMF}
NIST Digital Library of Mathematical Functions,
Sections 5.7 (Gamma series), 5.11 (asymptotics), 5.12 (beta integrals),
5.15 (polygamma functions), and 25.11 (Hurwitz zeta function).
\url{https://dlmf.nist.gov/5.7},
\url{https://dlmf.nist.gov/5.11},
\url{https://dlmf.nist.gov/5.12},
\url{https://dlmf.nist.gov/5.15},
\url{https://dlmf.nist.gov/25.11}.
Accessed 11 October 2026.

\end{thebibliography}

\end{document}
